\documentclass[11pt,a4paper]{article}
\usepackage[T1]{fontenc}
\usepackage{lmodern,microtype}
\usepackage[margin=25mm]{geometry}
\usepackage{amsmath,amssymb,amsthm,mathtools}
\usepackage{booktabs,longtable,array,enumitem,listings,xcolor,fancyhdr}
\usepackage[colorlinks=true,linkcolor=blue!45!black,citecolor=blue!45!black,urlcolor=blue!45!black,hypertexnames=false]{hyperref}
\setlength{\emergencystretch}{3em}
\setlength{\parskip}{.25em}
\setlength{\headheight}{14pt}
\pagestyle{fancy}\fancyhf{}
\fancyhead[L]{\small Exact local potentials and quantum path families}
\fancyhead[R]{\small Research manuscript}\fancyfoot[C]{\thepage}
\newtheorem{theorem}{Theorem}[section]
\newtheorem{lemma}[theorem]{Lemma}
\newtheorem{proposition}[theorem]{Proposition}
\newtheorem{corollary}[theorem]{Corollary}
\theoremstyle{definition}\newtheorem{definition}[theorem]{Definition}
\theoremstyle{remark}\newtheorem{remark}[theorem]{Remark}
\DeclareMathOperator{\rad}{rad}
\DeclareMathOperator{\diam}{diam}
\DeclareMathOperator{\ecc}{ecc}
\newcommand{\Z}{\mathbb Z}\newcommand{\R}{\mathbb R}\newcommand{\Q}{\mathbb Q}
\lstset{basicstyle=\ttfamily\scriptsize,breaklines=true,columns=fullflexible,keepspaces=true,showstringspaces=false,frame=single,rulecolor=\color{black!20},aboveskip=8pt,belowskip=8pt}
\title{Exact Local Potentials for the Randi\'c--Radius Problem\\[5pt]
\large The degree-seven case and new randomized quantum mixing families}
\author{Keio Research Notes\\\small With internal proof audits and exact computational certificates}
\date{2 November 2026}
\hypersetup{pdftitle={Exact Local Potentials for the Randic-Radius Problem},pdfauthor={Keio Research Notes}}
\begin{document}
\maketitle
\begin{abstract}
We prove that every finite connected simple graph of maximum degree at
most seven satisfies
\(\sum_{uv\in E(G)}1/\max(d_u,d_v)\ge\rho(G)\), where
\(\rho\) is the minimum eccentricity over vertex and edge-interior
centres with demand at the original vertices. Consequently its
Randi\'c index is at least \(\rad(G)-1/2\). The proof combines a
minimal-counterexample reduction with exact rational local potentials.
A first certificate bounds a hypothetical core by twelve vertices;
fourteen further dual certificates and one simplicity constraint reduce
this to eleven. An exhaustive integer-arithmetic core check completes
the proof. The source, certificate data, and verifiers are included in
the single TeX file.
Independently, we construct bounded compactly supported observation-time
densities giving exact uniform quantum mixing on \(P_{3p-1}\) for
every prime \(p\notin\{5,7\}\); the excluded cases are impossible
by an existing obstruction. A second construction covers
\(P_{2^A p^k-1}\) for primes \(p\ge5\), every \(A\ge1\), and
every \(k\ge2\).
These are advances on specified infinite classes. Neither the unrestricted
Fajtlowicz conjecture nor the classification of all quantum paths is
settled here. Historical priority and comparative research value are
not certified.
\end{abstract}

\section{Original targets and the exact scope of the results}
Fajtlowicz's 1988 conjecture asks whether every connected graph satisfies
\(R(G)\ge\rad(G)-1\), where
\[
 R(G)=\sum_{uv\in E(G)}\frac1{\sqrt{d_ud_v}}.
\]
We retain this original statement. Our graph theorem proves a stronger
inequality on the class \(\Delta(G)\le7\), with no restriction on the
number of vertices. It is a bounded-degree case, not a proof of the
unrestricted named conjecture. Cygan, Pilipczuk and \v{S}krekovski proved
the corresponding classical radius conclusions for maximum degree four
\cite{CPS}. Doig's later work treats cactus graphs and gives further
context \cite{Doig2023}.

For quantum walks, the target is exact uniformity of all averaged
transition probabilities under a single freely chosen observation-time
law. This is the randomized observation problem discussed by
Coutinho and Guo \cite{CoutinhoGuo2024} and studied for paths by Kang
\cite{Kang2026}. Our family classifications use this same notion.
They do not concern uniformity of the standard Ces\`aro time average,
instantaneous uniform mixing, or the separate rank-two average-mixing
problem. Published cases and obstructions are distinguished from the
new positive constructions below.

\paragraph{Starting material.}
We inspected the 102 numbered research drafts in the repository
\cite{Corpus}, at commit
\begin{center}\small\texttt{ecd87de3348dcb706a76bb12fde722d7a348a012}.\end{center}
Its timestamp is 2 November 2026, 19:38:17, UTC+09:00; a fresh fetch
confirmed this snapshot. The repository's own provenance statement
identifies it as an unrefereed collection and warns that it may contain
errors. We use its statements as research leads and reproduce the
analytic reduction needed from [087]. The preceding research note's
factorial core bound and odd-prime-power mixing theorem are starting
points, not results counted again in this round.

\paragraph{Methodological contribution and limits.}
The auxiliary graph problem is to assign rational potentials to local
degree types so that every permitted edge pays for the local vertex
inequality. Solving this finite problem converts a qualitative
minimal-counterexample reduction into a small, exact certificate.
The quantum construction instead prescribes moments on the actual
frequency-relation torus and then corrects them on a finite interval of
observation times. A positive quadratic density supplies correlations
with a dependent phase; a separate root-label construction handles a
fixed zero mode. These are concrete constructions relative to the
inspected source material. Linear programming duality, character
orthogonality and harmonic analysis are established tools; we do not
assert a historically unprecedented foundational principle.

The requested conjunction of an unchanged previously open named problem,
a complete unrestricted solution, historical novelty, and higher value
than all previous work has \emph{not} been established. Complete proofs
for the stated classes, a bounded literature search, and global priority
are different claims. Exact computations and independent audits support
the results here; they are not a proof-assistant formalization or external
peer review.


\section{The modified Randi\'c--radius inequality through degree seven}
\label{sec:g7-main}

For a finite simple graph put
\[
 M(G)=\sum_{uv\in E(G)}\frac1{\max(d_G(u),d_G(v))},
 \qquad R(G)=\sum_{uv\in E(G)}\frac1{\sqrt{d_G(u)d_G(v)}}.
\]
Let $\Gamma(G)$ replace every edge by a unit segment, and define the
\emph{vertex-demand absolute radius}
\[
 \rho(G)=\min_{x\in\Gamma(G)}\max_{v\in V(G)}d(x,v).
\]
The maximum here is over graph vertices, not over all points of the
metric realization. This distinction matters, for example on cycles.

\begin{theorem}\label{thm:g7-main}
Every connected finite simple graph of maximum degree at most seven satisfies
\[
 M(G)\ge\rho(G)\ge\operatorname{rad}(G)-\tfrac12.
\]
Consequently $R(G)\ge\operatorname{rad}(G)-\tfrac12$, and in particular
Fajtlowicz's inequality $R(G)\ge\operatorname{rad}(G)-1$ holds throughout
this class.
\end{theorem}

This is a bounded-degree result. It does not prove either the unrestricted
Fajtlowicz conjecture or the unrestricted Caporossi--Hansen conjecture.
For the classical Randi\'c--radius problem, the established
chemical-graph case has maximum degree four \cite{CPS}; corpus manuscript~\cite{Corpus087} develops the absolute-radius
reduction and proves the degree-five case. We reproduce the necessary
analytic reduction below. The new step is an exact degree-type potential:
it bounds every residual core by twelve vertices, and a second finite
certificate excludes order twelve. All rational certificates and their
standard-library verifier are included in the accompanying single source.
No claim that a literature search establishes worldwide priority is made.

\subsection{Elementary metric and deletion identities}

The minimum defining $\rho$ is attained at a vertex or an edge midpoint.
Indeed, on an edge $uv$, at coordinate $t\in[0,1]$, distance to a vertex
$w$ is $\min\{t+d(u,w),1-t+d(v,w)\}$; the two integer endpoint distances
differ by at most one. The upper envelope has a minimum at $0$, $1/2$,
or $1$. Thus
\begin{equation}\label{eq:g7-rhoformula}
 2\rho(G)=\min\left\{
 2\operatorname{rad}(G),\
 \min_{uv\in E(G)}\left(1+2\max_w\min\{d(u,w),d(v,w)\}\right)
 \right\},
\end{equation}
so $\rho\in\{\operatorname{rad}-1/2,\operatorname{rad}\}$.
Deleting a non-bridge edge cannot decrease $\rho$.

For a leaf $\ell$ with support $p$, $d=d_G(p)\ge2$, let $\lambda$ count
neighbors of $p$ other than $\ell$ whose full degrees are less than $d$.
Directly recomputing affected weights gives
\begin{equation}\label{eq:g7-leaf}
 M(G)-M(G-\ell)=\frac1d-\frac\lambda{d(d-1)}\ge0.
\end{equation}
For an edge $uv$ with both degrees at least two, let $\beta_u$ count
lower-degree neighbors of $u$ other than $v$, and define $\beta_v$
symmetrically. Then
\begin{equation}\label{eq:g7-edge}
 M(G-uv)-M(G)=
 \frac{\beta_u}{d_u(d_u-1)}+
 \frac{\beta_v}{d_v(d_v-1)}-\frac1{\max(d_u,d_v)}.
\end{equation}
An edge is called \emph{supportive} when this change is positive.

\begin{lemma}\label{lem:g7-leafmetric}
For a connected graph with at least three vertices and a leaf $\ell$,
\[
 \rho(G)-\tfrac12\le\rho(G-\ell)\le\rho(G).
\]
Deleting any set of leaves lowers $\rho$ by at most one. If two leaves
have the same support and $G$ is not a star, deleting one of them leaves
$\rho$ unchanged.
\end{lemma}
\begin{proof}
Distances between retained vertices and edge points do not change.
A center lying in a deleted pendant edge can be moved to its support,
proving the upper bound. Let $y$ be a vertex or midpoint center of
$G-\ell$, with value $\rho'$, and let $p$ support $\ell$. If
$d(y,p)\le\rho'-1/2$, the same point has value at most $\rho'+1/2$ in
$G$. Otherwise $d(y,p)=\rho'$, since both quantities are half-integers.
Move $y$ by $1/2$ toward $p$. All retained vertices are then within
$\rho'+1/2$, as is $\ell$. This proves the lower bound. For a set of
leaves, a center of the retained graph is within one extra unit of
every removed leaf. Finally, for two leaves at $p$, choose a center
of the graph after one deletion outside the remaining pendant edge:
if necessary, its projection to $p$ is at least as good. The remaining
leaf then has the same distance as the deleted one, proving equality.
\end{proof}

\begin{lemma}[Cut-vertex padding]\label{lem:g7-padding}
Suppose $v$ is a cut vertex and $V(G)\setminus\{v\}=A\sqcup B$ with no
edge from $A$ to $B$, where both $G[A]$ and $G[B]$ contain an edge.
Let $H_A=G[A\cup\{v\}]$, and obtain $G_A$ by adding $d_B(v)$ new leaves
at $v$. Define $G_B$ symmetrically. Then degrees of retained vertices
are preserved, and
\[
 M(G)=M(G_A)+M(G_B)-1,
 \qquad \rho(G)\le\rho(G_A)+\rho(G_B)-1.
\]
\end{lemma}
\begin{proof}
The two padded graphs preserve all original edge weights, while their
new leaf weights sum to $(d_A(v)+d_B(v))/d_G(v)=1$.
Put $\rho_A=\rho(G_A)$ and choose a center $y_A\in\Gamma(H_A)$;
a center in a new pendant edge can be projected to $v$. Set
$s_A=d(y_A,v)$. Since $G_A$ has a leaf at $v$, $s_A\le\rho_A-1$.
Define $y_B,s_B$ similarly. On a geodesic from $y_A$ through $v$ to
$y_B$, with signed coordinate $p\in[-s_A,s_B]$, every vertex is within
\[
 \max\{\rho_A+s_A+p,\rho_B+s_B-p\}.
\]
If the balancing point lies in this interval, the minimum of this
upper bound is $(\rho_A+\rho_B+s_A+s_B)/2\le\rho_A+\rho_B-1$.
If it lies outside, one endpoint gives at most $\max(\rho_A,\rho_B)$,
which is also at most $\rho_A+\rho_B-1$ because both radii are at least
one.
\end{proof}

\subsection{The residual core and the source budget}

\begin{proposition}\label{prop:g7-reduction}
If $M\ge\rho$ fails among connected graphs of maximum degree at most
$D\ge3$, a counterexample exists of the form $G=K+L$, where $K$ is
$2$-connected of order at least three and $L$ consists of at most one
leaf per core vertex, such that:
\begin{enumerate}[label=(\roman*)]
\item all full core degrees lie in $[3,D]$, every core edge is supportive,
and every core vertex has a lower-degree neighbor in $G$;
\item $\delta:=M(G)-M(K)<1$;
\item exactly $s\le2$ core vertices have core degree two. They have full
degree three, carry leaves, are independent, and have core neighbors
of full degree at least four.
\end{enumerate}
\end{proposition}
\begin{proof}
Choose a counterexample minimizing lexicographically its cyclomatic
number $|E|-|V|+1$ and then its number of edges. Every non-bridge edge
is supportive; otherwise deleting it decreases cyclomatic number,
does not increase $M$, and does not decrease $\rho$.

At a cut vertex, at most one component after deletion can contain an
edge. Otherwise apply Lemma~\ref{lem:g7-padding}. Each padded graph
has the same degree bound and smaller lexicographic size: its
cyclomatic number cannot increase, and if unchanged, the omitted
component's internal edge makes the edge count smaller. Minimality
and the two padding inequalities contradict $M(G)<\rho(G)$.

No leaf can have a degree-two support, except in the already valid
star $P_3$: deleting such a leaf drops $M$ by exactly $1/2$ and drops
$\rho$ by at most $1/2$. No vertex carries two leaves, because deleting
one does not increase $M$ and, outside stars, leaves $\rho$ unchanged.
Stars satisfy $M\ge\rho$. If $G$ were a tree, the cut-vertex property
would make it a star or double star; the preceding restrictions leave
only $P_4$ among the latter, already excluded by the degree-two support
argument. Thus $G$ contains a $2$-connected block $K$ of order at least
three. The cut-vertex property forces every vertex outside $K$ to be
an attached leaf, proving the asserted form.

A core vertex $u$ with no lower-degree neighbor would have no leaf.
For a neighbor $v$ one has $d_v\ge d_u$, and
\eqref{eq:g7-edge}, with $\beta_u=0$ and $\beta_v\le d_v-1$, contradicts
supportiveness. A full core degree of two is consequently impossible.
A full degree-three vertex has to carry a leaf, since all core degrees
are at least three; two such vertices cannot be adjacent because
\eqref{eq:g7-edge} would then give zero. Conversely a core-degree-two
vertex must carry its one leaf and have full degree three.

Deleting any nonempty subset $L'$ of leaves produces a smaller graph,
so minimality and Lemma~\ref{lem:g7-leafmetric} give
$M(G)-M(G-L')<\rho(G)-\rho(G-L')\le1$. In particular $\delta<1$.
If $L'$ consists of the leaves at the full degree-three vertices, its
exact drop is $s/3$: their other neighbors retain degrees at least
four. Hence $s/3<1$, so $s\le2$.
\end{proof}

A \emph{source} is a core vertex without a lower-degree core neighbor.
Every source has a leaf. Sources are independent: unequal degrees
contradict the definition at the higher endpoint, and for equal degree
$d\ge3$ the deletion change is $(3-d)/[d(d-1)]\le0$.
Deleting just their leaves changes no retained edge weight. Therefore
\begin{equation}\label{eq:g7-source}
 \sum_{v\text{ a source}}\frac1{d_v}\le\delta<1.
\end{equation}
In particular, if $N_d$ counts full core degrees and $d_0$ is the least
occupied degree, then $N_{d_0}<d_0$. This budget will remove degree
profiles before the final finite certificate.

\subsection{Degree-type potentials and the order-twelve exclusion}

In this subsection, a decorated core has minimum core degree at least
two and at most one added leaf at each core vertex; it is supportive
if every core edge is supportive. For a core vertex define its type
$t=(d,p,l)$ by its full degree, leaf indicator, and number of
lower-degree neighbors in $G$. Put $g_d=d(d-1)$.
For $D=7$ all possible types belong to the following set of size $45$:
\[
 \mathcal T=\{(3,1,1)\}\ \cup\
 \{(d,p,l):4\le d\le7,\ p\in\{0,1\},\ 1\le l\le d\}.
\]
Including types that cannot actually occur only enlarges our relaxation.
For an edge joining $t=(d,p,l)$ to $u=(e,q,m)$ with $d<e$,
supportiveness is equivalent to
\begin{equation}\label{eq:g7-compatible}
 \frac{l}{g_d}>\frac{e-m}{g_e};
 \qquad\text{for }d=e\text{ it is equivalent to }l+m\ge d.
\end{equation}
These tests leave $802$ allowed unordered type pairs, including pairs
of the same type. A vertex of type $t$ has $l-p$ lower-degree core
neighbors and $d-l$ equal-or-higher-degree core neighbors.

Simultaneously deleting all leaves gives the exact identity
\begin{equation}\label{eq:g7-delta}
 \delta=\sum_t n_t\frac{p(d-l)}{g_d}
 -\sum_{\substack{uv\in E(K)\\d_u=d_v=d,\ p_u=p_v=1}}\frac1{g_d}.
\end{equation}
Indeed, deleting a host's leaf removes $1/d$ and increases each edge
to a lower-degree retained neighbor by $1/g_d$. Equal-degree edges
increase only when both endpoints are hosts, producing the final
correction.

\begin{lemma}[Rational potential certificate]\label{lem:g7-potential}
Every supportive decorated core with full degrees in $[3,7]$ satisfies
\[
 |V(K)|\le\frac{875}{68}\,\delta.
\]
\end{lemma}
\begin{proof}
The included certificate assigns rational numbers $\alpha_t,\beta_t$
to the $45$ types. With $C=875/68$, it verifies exactly
\begin{align}
 (l-p)\alpha_t+(d-l)\beta_t+C\frac{p(d-l)}{g_d}&\ge1,
                                                    \label{eq:g7-vcert}\\
 \beta_t+\alpha_u&\le0 &&(d<e),\label{eq:g7-ecert1}\\
 \beta_t+\beta_u+C\frac{pq}{g_d}&\le0 &&(d=e),\label{eq:g7-ecert2}
\end{align}
where the last two inequalities range over allowed pairs only.
These are $45+802$ explicit rational comparisons. Sum
\eqref{eq:g7-vcert} over vertices. Each unequal-degree edge contributes
$\beta$ at its lower endpoint and $\alpha$ at its higher endpoint;
an equal-degree edge contributes two $\beta$ terms. Applying
\eqref{eq:g7-ecert1}--\eqref{eq:g7-ecert2} and
\eqref{eq:g7-delta} proves the lemma. The supplied verifier checks the
certificate with integer fractions and does not invoke an optimizer.
\end{proof}

Thus $\delta<1$ forces $|K|\le12$. To exclude twelve, let
$N=(N_3,N_4,N_5,N_6,N_7)$. The exact candidate list consists of
\begin{equation}\label{eq:g7-profiles}
 N_d\ge0,\quad\sum N_d=12,\quad N_3\le2,
 \quad N_{d_0}<d_0\text{ for the least occupied degree }d_0.
\end{equation}
There are $896$ such integer profiles. At a fixed profile, discard
any type failing
\begin{equation}\label{eq:g7-typefilter}
 N_d>0,\qquad l-p\le\sum_{e<d}N_e,
 \qquad d-l\le\sum_{e\ge d}N_e-1.
\end{equation}
Every discarded type violates a necessary distinct-neighbor count.

\begin{lemma}[Finite profile exclusion]\label{lem:g7-profiles}
No supportive decorated core of order twelve, with full degrees in
$[3,7]$ and $\delta<1$, exists.
\end{lemma}
\begin{proof}
We specify the finite certificate mathematically so that its checks
are independent of how its coefficients were found. Let $x_{tu}$ be
the number of allowed core edges of type pair $tu$. The lower and
upper incidence equations at type $t=(d,p,l)$ are
\[
 \sum_{u:d_u<d}x_{tu}=(l-p)n_t,
 \qquad
 \sum_{u:d_u\ge d}x_{tu}=(d-l)n_t,
\]
where a same-type edge contributes twice to the second sum. Add the
five degree-count equations $\sum_{t:d_t=d}n_t=N_d$.
The objective is the right side of~\eqref{eq:g7-delta}.

A dual vector comprises two unrestricted rational multipliers per
type and five degree-count multipliers. It is a certificate of a
lower bound on $\delta$ if its coefficient on every retained vertex
or edge variable is at most that variable's coefficient in
\eqref{eq:g7-delta}. The resulting bound is the dot product of its
last five entries with $N$. All variables being nonnegative, this
implication follows by summing inequalities, without linear-programming
duality as an external assumption.

The included data contain $14$ such rational vectors. The verifier
reconstructs all profiles~\eqref{eq:g7-profiles}, retains the columns
specified by~\eqref{eq:g7-typefilter} and~\eqref{eq:g7-compatible},
and confirms that at least one vector proves $\delta\ge1$ for each
of $895$ profiles. Its sole residual profile is
\[
 N=(0,3,3,3,3).
\]
For this profile let $Z$ be the degree-six vertices with $l=1$.
They form an independent set by~\eqref{eq:g7-compatible}. Every one
needs five distinct equal-or-higher-degree neighbors. There are only
six vertices of degree at least six, so $|Z|\le1$ (the case $Z$ empty
is immediate). Incorporating this one additional inequality, the
residual rational certificate proves
\[
 \delta\ge\frac{103}{84}>1.
\]
Its multiplier on $|Z|\le1$ is $-61/420$; all other added capacity
multipliers are zero. The verifier checks every coefficient and the
exact lower bound. This excludes the residual case and proves the
lemma.
\end{proof}

The certificate data are small: one $45$-type potential, fourteen
$95$-entry profile vectors, and one residual vector. The single source
contains their compressed numerical payload and a Python verifier
using only \texttt{fractions}, \texttt{json}, and other standard-library
modules. The reconstructed profile count, uncovered profile, and exact
residual bound are all asserted by the verifier. No floating-point
solver result is used as a proof certificate.

\subsection{The finite core inequality and conclusion}

\begin{lemma}[Exhaustive finite core certificate]\label{lem:g7-smallcores}
If $K$ is $2$-connected, has at most eleven vertices, maximum degree
at most seven, minimum degree at least two, and exactly $s\le2$
vertices of degree two, then
\begin{equation}\label{eq:g7-coreineq}
 \rho(K)\le M(K)-1+\frac s3.
\end{equation}
\end{lemma}
\begin{proof}
Generate every $2$-connected graph of the stated order and degree
bounds with \texttt{geng}; retain precisely those with $s\le2$.
For each retained graph, compute all distances by breadth-first
search and $2\rho$ by~\eqref{eq:g7-rhoformula}. Let
$W=\operatorname{lcm}(6,2,3,\ldots,\Delta(K))$ and let $J=WM(K)$,
computed as a sum of integers $W/\max(d_u,d_v)$. The exact integer
comparison is
\[
 3W(2\rho)\le6J-6W+2Ws.
\]
No violation occurs. The enumeration and integer checker are included
in the reproducibility material. The counts below distinguish the
unrestricted-degree checks through order nine from the degree-seven
checks at orders ten and eleven; the former are a safe enlargement.
\begin{center}
\begin{tabular}{rrrr}
\toprule
Order&Maximum-degree filter&Eligible cores&Minimum slack\\
\midrule
4&none&2&$0$\\
5&none&6&$1/6$\\
6&none&44&$1/6$\\
7&none&385&$1/4$\\
8&none&$6\,268$&$3/10$\\
9&none&$179\,303$&$1/3$\\
10&7&$7\,355\,341$&$5/14$\\
11&7&$473\,655\,792$&$1/2$\\
\bottomrule
\end{tabular}
\end{center}
Slack means $M(K)-1+s/3-\rho(K)$. Order three has no eligible graph.
An independent Python implementation, using \texttt{Fraction} and
separate graph decoding and distance calculations, reproduced every
count and slack through order nine, for $186\,008$ eligible cores.
The higher-order C run was newly executed, rather than inferred from
the earlier corpus table. Enumeration completeness uses the standard
\texttt{nauty} generator; this is a computer-assisted finite proof,
not a proof-assistant formalization.
\end{proof}

\begin{proof}[Proof of Theorem~\ref{thm:g7-main}]
If a counterexample of maximum degree seven exists,
Proposition~\ref{prop:g7-reduction} supplies $G=K+L$ with $\delta<1$.
Lemmas~\ref{lem:g7-potential} and~\ref{lem:g7-profiles} give $|K|\le11$.
Let $L_3$ be the leaves attached at the $s$ full degree-three core
vertices. Their exact deletion drop is $s/3$, and subsequent leaf
deletions do not increase $M$. Thus $M(K)\le M(G)-s/3$.
Adding leaves raises $\rho$ by at most one. By
Lemma~\ref{lem:g7-smallcores},
\[
 \rho(G)\le\rho(K)+1\le M(K)+s/3\le M(G),
\]
a contradiction. Finally, $R\ge M$ termwise and
$\rho\ge\operatorname{rad}-1/2$ by~\eqref{eq:g7-rhoformula}.
\end{proof}

For reproducibility, the generator was built from the
\texttt{lonnen/nauty} source mirror, release marker
\texttt{This\_is\_nauty2\_9\_0.txt}, commit
\texttt{3915d891af81630853ab485f2696828a31253cd9}.
The high-order commands are
\begin{lstlisting}[language=bash,basicstyle=\ttfamily\footnotesize,breaklines=true]
geng -C -d2 -D7 -q 10 | graph_corechk
geng -C -d2 -D7 -q 11 | graph_corechk
python3 graph_verify_compact.py graph_degree7_compact_certificate.json
\end{lstlisting}
The generator, independent rational certificate, and two graph checkers
serve different roles: generation establishes finite coverage,
the rational certificate proves the reduction in order, and the
checkers evaluate the finite core inequality. The bounded-degree
restriction remains essential to the result stated here.


\section{Randomized uniform mixing when the path modulus is three times a prime}
\label{sec:three-prime-paths}

For a graph with adjacency matrix $A$, a probability law for the observation
time $T\geq0$ gives \emph{randomized uniform mixing} if
\[
 \mathbb E\bigl|(e^{iTA})_{uv}\bigr|^2=\frac1{|V|}
 \qquad(u,v\in V).
\]
The same law must work for every initial vertex. This is the ordinary
adjacency quantum walk with a freely chosen time law. It is different from
the standard Ces\`aro average, which fixes the averaging protocol to
$\lim_{L\to\infty}L^{-1}\int_0^L$. In particular, nothing here settles
the rank-two Ces\`aro question in \cite[Open Problem~4.1]{CoutinhoGuo2024}.
The randomized path question is discussed separately in
\cite[Section~4.2]{CoutinhoGuo2024}.

\begin{theorem}[The three-times-prime family]
\label{thm:three-prime-paths}
Let $p$ be a prime. The unweighted path $P_{3p-1}$ admits randomized
uniform mixing if and only if $p\notin\{5,7\}$. In every positive case,
the time law can be chosen to have a bounded, compactly supported
probability density on $[0,\infty)$.
\end{theorem}

Kang's negative theorem already excludes $p=5,7$, and his positive
theorem includes $p=2$ \cite{Kang2026}. The additional cases established
below are the surviving members of this particular arithmetic family.
This is not an all-path classification.

\subsection{Spectral target and exact realization in time}
Write $q=n+1$ for the modulus of $P_n$. Its eigenvalues, normalized
eigenvectors and spectral projectors are
\[
 \lambda_j=2\cos(\pi j/q),\qquad
 v_j(u)=\sqrt{2/q}\sin(\pi ju/q),\qquad
 E_j=v_jv_j^{\mathsf T}\quad(1\leq j,u<q).
\]
Let $R$ reverse the vertex order and let $\circ$ denote the entrywise
product. Two useful identities are
\begin{equation}\label{eq:new-path-spectral}
 D_0:=\sum_j E_j\circ E_j
 =\frac1q\left(J+\frac{I+R}{2}\right),\qquad
 \sum_{j\equiv\ell\pmod2}E_j\circ E_\ell=\frac{I+R}{2}.
\end{equation}
For the first identity, expand the product of two squared sines and use
\(\sum_{j=1}^{q-1}\cos(2\pi aj/q)=q\mathbf1_{q\mid a}-1\).
For the second, use $Rv_j=(-1)^{j+1}v_j$ and expand
\[
 \frac12\left[
 \left(\sum_jE_j\right)\circ\left(\sum_jE_j\right)
 +\left(\sum_j(-1)^jE_j\right)\circ
  \left(\sum_j(-1)^jE_j\right)\right].
\]
Consequently, phase coherences equal to $1$ on the diagonal, $-c$ on
distinct equal-parity pairs, and $0$ on opposite-parity pairs give the
mixing matrix
\begin{equation}\label{eq:new-path-target}
 (1+c)D_0-c\frac{I+R}{2}=\frac Jn\qquad\text{when }c=1/n.
\end{equation}
Only real parts of these coherences enter the mixing matrix, although
our constructions satisfy the displayed complex moments as well.

When $q$ is odd, each pair $j,q-j$ has one odd and one even member,
and $\lambda_{q-j}=-\lambda_j$. Thus it suffices to construct the
$m=(q-1)/2$ odd-index phases with
\begin{equation}\label{eq:new-odd-moments}
 \mathbb E W_i\overline W_j=-c\quad(i\ne j),\qquad
 \mathbb E W_iW_j=0\quad(\text{all }i,j),
\end{equation}
and take even phases to be their reflected conjugates. All these laws
must be supported on the actual orbit closure of the time flow.

The next lemma supplies the passage from that closure to an actual
observation-time law. It will also be used for other path families.

\begin{lemma}[Exact finite-time realization]\label{lem:exact-finite-time}
Let $\mathcal T$ be a finite-dimensional torus, and let
$\gamma(t)$, $t\geq0$, be an equidistributed linear flow in $\mathcal T$.
Let $F_1,\ldots,F_D$ be distinct nontrivial cosine characters, identified
modulo sign. Suppose a probability measure $\mu$ on $\mathcal T$ has the
desired finite moment vector $b_j=\int F_j\,d\mu$ and satisfies
$\mu\geq\eta H$ for some $\eta>0$, where $H$ is Haar probability.
Then a bounded, compactly supported probability density $\rho$ on
$[0,\infty)$ satisfies
$\int_0^\infty F_j(\gamma(t))\rho(t)\,dt=b_j$ for every $j$.
\end{lemma}
\begin{proof}
Inversion-symmetrize $\mu$, which preserves its cosine moments and its
Haar lower bound. Convolve it with a nonnegative product Fej\'er
approximate identity. The resulting trigonometric-polynomial density
is at least $\eta$, and its finitely many specified moments converge to
$b_j$. Distinct cosine characters have Haar inner products
$\int F_iF_j\,dH=\delta_{ij}/2$. If the moment error is $\delta_j$,
adding $2\sum_j\delta_jF_j$ corrects every moment, preserves the
integral, and changes the sup norm by at most $2\sum_j|\delta_j|$.
Choosing the approximation sufficiently close gives a trigonometric
polynomial $f\geq\eta/2$ with exactly the desired moments.

Put $F_0=1$, $b_0=1$, and regard $F$ and $b$ as column vectors.
Along the flow define
\[
 G_L=\frac1L\int_0^L F(\gamma(t))F(\gamma(t))^{\mathsf T}\,dt,
 \qquad b_L=\frac1L\int_0^L F(\gamma(t))f(\gamma(t))\,dt.
\]
Equidistribution gives
$G_L\to\operatorname{diag}(1,1/2,\ldots,1/2)$ and $b_L\to b$.
Therefore, for all sufficiently large finite $L$,
$a_L=G_L^{-1}(b-b_L)$ exists and tends to zero. Set
\[
 \rho(t)=\frac{f(\gamma(t))+a_L^{\mathsf T}F(\gamma(t))}{L}
             \mathbf1_{[0,L]}(t).
\]
For sufficiently large $L$ this is nonnegative, since $f\geq\eta/2$.
The defining linear system gives $\int\rho F=b$, including
$\int\rho=1$. Boundedness and compact support are immediate.
\end{proof}

For the path constructions, a law satisfying \eqref{eq:new-odd-moments}
at some $c>c_0=1/n$ supplies the lemma's Haar component by taking
\begin{equation}\label{eq:new-haar-mix}
 \mu=(1-c_0/c)H+(c_0/c)\nu_c.
\end{equation}
Every relevant off-diagonal Hermitian or holomorphic character has
zero Haar integral: distinct odd eigenvalues are neither equal nor
opposite, and none is zero. The mixture therefore has exactly the
target moments. For the full spectrum, the analogous assertion uses
the nonzero eigenvalue differences. Repeated cosine characters are
listed just once when applying Lemma~\ref{lem:exact-finite-time}.

\subsection{The exact relation torus for \texorpdfstring{$q=3p$}{q=3p}}
Assume $p\geq5$, and put
\[
 r=\frac{p-1}{2},\qquad m=3r+1,\qquad n=6r+2.
\]
The odd indices partition into the singleton $p$ and the $r$ triples
\[
 \{j,2p-j,2p+j\},\qquad j=1,3,\ldots,p-2.
\]
Each triple of frequencies sums to zero. Exactly one index in each
triple is divisible by $3$; call its phase $U_a$. The sum of those
frequencies is
\begin{equation}\label{eq:new-global-relation}
 \sum_{b=1,3,\ldots,p-2}2\cos(\pi b/p)=1=\lambda_p.
\end{equation}
The trigonometric sum is $\sin((p-1)\pi/p)/\sin(\pi/p)=1$.
Hence the singleton phase satisfies $W_0=\prod_{a=1}^r U_a$.

These $r+1$ relations generate the entire integer relation lattice.
Indeed, the real cyclotomic field $K=\mathbb Q(2\cos(\pi/q))$ has
degree $\varphi(6p)/2=p-1=2r$. The odd-index eigenvalues span $K$:
powers of $\lambda_1$ expand into $1$ and frequencies
$2\cos(a\pi/q)$, periodicity and reflection reduce all these to the
odd modes, and $1=\lambda_p$ is already present. Thus their rational
relation space has dimension $m-2r=r+1$. The disjoint triple relations
and \eqref{eq:new-global-relation} are independent, because only the
latter has a nonzero singleton coordinate. To check saturation, take
an integral vector in their rational span. Its singleton coordinate
forces the coefficient of the extra relation to be integral. After
subtracting that multiple, the three coordinates in each disjoint
triple force its remaining coefficient to be integral.

The character criterion for a linear torus flow now shows that its
closure is exactly $\mathcal T\cong(S^1)^{2r}$ with characters
\begin{equation}\label{eq:new-three-torus}
 \chi_{a1}=u_a,\qquad \chi_{a2}=v_a,\qquad
 \chi_{a3}=-u_a-v_a,\qquad
 \gamma=\sum_{a=1}^r u_a.
\end{equation}
Here $u_a,v_a$ are a free integral character basis, $U_a$ is the phase
of $u_a$, and $W_0$ is the phase of $\gamma$.
Every nontrivial character of this torus has a nonzero time frequency;
its time average tends to zero. This also proves the equidistribution
needed in Lemma~\ref{lem:exact-finite-time}.

\subsection{Anticorrelating the triple modes while cancelling singleton moments}
Call the $3r$ triple phases the old modes. For $r\geq4$, we first
construct a law $\nu_d$ with old-mode off-diagonal Hermitian moments
$-d$, all holomorphic moments zero, and all singleton--old Hermitian
moments zero. The allowable interval is
\begin{equation}\label{eq:new-dmax}
 0<d\leq d_{\max}:=\frac{r-2}{3r^2-r-4}.
\end{equation}
Put $B=\sqrt{3-6d}$ and choose $\theta$ with
$1+2\cos\theta=B$. For each triple use an independently uniform
permutation of
\begin{equation}\label{eq:new-triple-template}
 Y_a(1,e^{i\theta},e^{-i\theta}),\qquad Y_a^3=1.
\end{equation}
The designated $U_a$ is its fixed indexed coordinate, and the product
of the three entries is one. It remains to construct exchangeable
labels with pair correlation
\[
 \mathbb E Y_a\overline Y_b=-A_d,
 \qquad A_d=\frac{3d}{1-2d}\quad(a\ne b).
\]

Choose positive $s_1+s_2=r$ with neither size divisible by $3$.
One can take $(1,r-1)$ if $r\equiv0,2\pmod3$, and $(2,r-2)$ if
$r\equiv1\pmod3$. Uniformly partition the $r$ positions into these
two sizes. Within each group distribute labels among the three cube
roots with counts differing by at most one, uniformly permute, and
apply an independent common uniform cube-root rotation. In a group
of size $s\not\equiv0\pmod3$, the label sum has squared modulus
one. Its ordered off-diagonal correlation sum is therefore $-(s-1)$.
The groups have zero cross correlations, so exchangeability makes the
overall off-diagonal label correlation
\[
 -A_{\max}=-\frac{r-2}{r(r-1)}.
\]
Mix this distribution with independent uniform labels to attain any
$0\leq A_d\leq A_{\max}$, which is precisely
\eqref{eq:new-dmax}. Both component laws may be generated with the
same two independent group rotations.

The within-triple Hermitian correlation in
\eqref{eq:new-triple-template} is $(B^2-3)/6=-d$.
Between different triples it is $-B^2A_d/9=-d$.
All old-mode holomorphic second moments vanish: use a group's common
rotation, of charge two if both factors are in that group and of
charge one otherwise.

The global relation $W_0=\prod U_a$ is now useful. Under the common
rotation $\omega$ of a group of size $s$, it changes by $\omega^s$.
For a singleton--old Hermitian or holomorphic moment, rotate the group
not containing the old phase. Since $3\nmid s$, the moment vanishes.
For $\mathbb E W_0^2$, either group has nonzero charge $2s$ modulo $3$,
so that moment also vanishes. This proves all the asserted properties
of $\nu_d$ on the exact torus \eqref{eq:new-three-torus}.

\subsection{A positive density supplies the missing singleton correlations}
Let $S$ be the sum of the $L=3r$ old phases. For $\alpha>0$ define
on Haar measure $H$ of the actual torus
\begin{equation}\label{eq:new-star-density}
 f_\alpha=\frac{|W_0-\alpha S|^2}{Z},\qquad Z=1+L\alpha^2.
\end{equation}
For $r\geq4$, its moments are
\begin{equation}\label{eq:new-star-moments}
 \mathbb E_{f_\alpha H}W_i\overline W_j=\frac{\alpha^2}{Z}
 \ (i\ne j\text{ old}),\qquad
 \mathbb E_{f_\alpha H}W_0\overline W_i=-\frac\alpha Z,
 \qquad \mathbb E_{f_\alpha H}W_iW_j=0
 \ \text{for all modes}.
\end{equation}
Here is the character check, including possible collisions. In each
block the old characters are $u,v,-u-v$; their six nonzero ordered
differences are distinct. Differences supported on two different
blocks are likewise distinguished by their two nonzero components.
Thus all nonzero ordered old differences are distinct. No old sum
(including a doubled character) equals an old difference: simultaneous
multiplication of every triple by a cube root gives charge two to a
sum and charge zero to a difference. No old characters are negatives
of one another.
The character $\gamma$ in \eqref{eq:new-three-torus} has support on
all $r$ blocks, and $\gamma\pm\chi_i$ has support on at least $r-1$.
For $r\geq4$ these cannot equal old sums or differences, supported on
at most two blocks. Moreover, $2\gamma$ cannot be a sum or difference
of two old characters. These observations rule out all unintended
collisions among the constant term, old differences, and
$\pm(\gamma-\chi_i)$ in the expansion of
\eqref{eq:new-star-density}. They also rule out every holomorphic
second-moment character from this Fourier support. Character
orthogonality therefore proves normalization and
\eqref{eq:new-star-moments}.

Mix $f_\alpha H$ and $\nu_d$ with weights $t$ and $1-t$. To obtain
the common Hermitian value $-c$, set
\begin{equation}\label{eq:new-mixture-parameters}
 t=cZ/\alpha,\qquad d=\frac{c(1+\alpha)}{1-t}.
\end{equation}
Then the singleton--old value is $-c$, the old--old value is
$-(1-t)d+c\alpha=-c$, and all holomorphic moments remain zero.

Take $\alpha=1/4$ and $c_0=1/(6r+2)$. At $c=c_0$ the parameters are
\[
 t_0=\frac{3r+16}{24r+8}<1,\qquad d_0=\frac5{21r-8}.
\]
For $r\geq7$ we have $d_0<d_{\max}$, since
\[
 (r-2)(21r-8)-5(3r^2-r-4)=6r^2-45r+36>0.
\]
Explicitly define
\begin{equation}\label{eq:new-explicit-c}
 c_{\max}=\frac{4(r-2)}{18r^2+5r-52},\qquad
 c=\frac{c_0+c_{\max}}2.
\end{equation}
The same inequality gives $c_{\max}>c_0$.
Substitution in \eqref{eq:new-mixture-parameters} shows $0<t<1$ and
$0<d<d_{\max}$: indeed $c_{\max}$ is exactly the value for
$d=d_{\max}$, and its denominator is positive before $t$ reaches one.
We have therefore constructed \eqref{eq:new-odd-moments} at a strict
value $c>c_0$. The Haar mixture \eqref{eq:new-haar-mix},
Lemma~\ref{lem:exact-finite-time}, and
\eqref{eq:new-path-target} prove the theorem for every prime $p\geq17$.

\subsection{Small positive cases and the two obstructions}
For $p=11,13$, the complete rational certificates in
Section~\ref{sec:exact-33-39} give laws on the actual phase torus at
$c=1/17>1/32$ and $c=1/20>1/38$, respectively. In each case the law
is explicitly symmetrized under complex conjugation; thus the exact
cosine-moment verification also establishes the stated complex
moments. The same Haar mixture and finite-time lemma apply.

For $p=3$, $q=9$, the odd frequencies at indices $1,3,5,7$ have the
single relation $\lambda_1+\lambda_5+\lambda_7=0$; the singleton is
$\lambda_3=1$. This is the entire integral relation lattice, by the
same field-degree and saturation argument (the field has degree three).
Use $c=1/6>1/8$, $B=\sqrt2$, and uniformly permute
$(1,e^{i\theta},e^{-i\theta})$ in the triple, with
$1+2\cos\theta=B$. Independently take the singleton in $\{1,-1\}$
with mean $-1/(2\sqrt2)$, and apply a common uniform cube-root rotation
to all four phases. The triple product remains one. The within-triple
correlation is $(B^2-3)/6=-1/6$, the singleton--triple correlation is
$[-1/(2\sqrt2)]B/3=-1/6$, and all holomorphic moments vanish under
the common rotation. Haar mixing and the lemma prove the required
bounded compactly supported density.

For $p=2$, $q=6$, the spectrum is $\sqrt3,1,0,-1,-\sqrt3$ and the
two positive frequencies are rationally independent. On its two-torus
take the phase ramp $z_j=e^{i(3-j)\xi}$, $1\leq j\leq5$, and give
$\xi\in\mathbb R/(2\pi\mathbb Z)$ the density
\[
 1-2c(\cos2\xi+\cos4\xi),\qquad c=9/40>1/5,
\]
relative to normalized Haar measure. This density is at least $1/10$,
and its coherences are $-c$ at nonzero even index differences and zero
at odd differences. It respects the spectral pairing and fixed zero
mode. Haar mixing on the actual two-torus and
Lemma~\ref{lem:exact-finite-time} finish this case; this is the small
instance of Kang's phase-ramp construction \cite{Kang2026}.

Finally, $p=5$ and $p=7$ give moduli $15$ and $21$, respectively.
Kang's forced-coherence theorem and parity-collision obstruction
\cite[Theorems S13 and S18]{Kang2026} exclude every time probability
law. Concretely, at $q=15$,
$\lambda_1-\lambda_3=\lambda_4-\lambda_5$, and at $q=21$,
$\lambda_3-\lambda_6=\lambda_7-\lambda_9$.
In either identity one pair has equal parity and the other opposite
parity, requiring the same cosine moment to equal both $-1/(q-1)$
and zero. This completes Theorem~\ref{thm:three-prime-paths}.

\paragraph{Scope and priority.}
The added proof combines two independently rotated label groups with
the nonnegative density \eqref{eq:new-star-density} to accommodate
the extra global frequency relation. Torus approximation, spectral
coherences and moment correction have prior antecedents; they are
not presented as unprecedented general principles. A targeted review
of the primary sources available through 2 November 2026 found no
earlier identical $q=3p$ classification. This is a bounded literature
assessment, not a proof of universal absence. The classification of
randomized uniform mixing for all paths remains unresolved here.


\subsection{Exact certificates at \texorpdfstring{$q=33$ and $q=39$}{q=33 and q=39}}
\label{sec:exact-33-39}
For completeness, the following finite certificates supply the two exceptional
positive cases used in the $q=3p$ theorem. They use only rational and integer
arithmetic, and all their data are included below.

Let $\alpha=2\cos(\pi/q)$ and let $P$ be its monic minimal polynomial of
degree $D=\varphi(2q)/2$. Each digit string $(z_0,\ldots,z_{D-1})$ assigns
$\exp(2\pi i z_a/3)$ to the phase of $\alpha^a$.
For an odd eigenvalue $\lambda_j=f_j(\alpha)$, with $f_j$ reduced modulo $P$,
the corresponding phase is
$W_j=\exp(2\pi i\sum_a[f_j]_a z_a/3)$.
Each string is given its displayed integer numerator divided by the common
denominator. Independently apply a uniform Galois permutation
$j\mapsto \min(aj\bmod 2q,\,2q-(aj\bmod 2q))$, with
$1\leq a<q$ and $\gcd(a,2q)=1$.
Finally symmetrize under complex conjugation.

The exact moments are
\[
 \mathbb E(W_i\overline W_j)=-c\quad(i\ne j),\qquad
 \mathbb E(W_iW_j)=0\quad(\hbox{all }i,j),
\]
with $c=1/17>1/32$ for $q=33$ and $c=1/20>1/38$ for $q=39$.
The weights are strictly positive. Mixing with Haar measure and applying the
finite-time realization lemma proves exact randomized uniform mixing for
$P_{32}$ and $P_{38}$.

The following self-contained Python program verifies every original moment,
not just a symmetry-reduced linear system. It uses only the Python standard
library. It independently reconstructs $\Phi_{2q}$, checks
$z^DP(z+z^{-1})=\Phi_{2q}(z)$, reconstructs all eigenvalue polynomials by
$f_0=2$, $f_1=x$, $f_j=xf_{j-1}-f_{j-2}$, and checks all moments using
$2\cos(2\pi a/3)\in\{2,-1\}$. The two successful outputs are exact integer
identities, with no floating-point tolerance.

{\footnotesize
\begin{verbatim}
"""Exact self-contained certificates for q=33 and q=39. Standard library only."""
import math
from fractions import Fraction
from functools import lru_cache

DATA = [
(33, '1/17', 93636,
 [1, 12, 12, -43, -43, 34, 34, -10, -10, 1, 1],
 '''
1000100222:861 1111112012:6966 2211020212:738 1020111021:3291
2021012210:9426 2110221112:492 1010122220:537 1110122002:279 0022122120:8934
0211221201:6355 2011120111:6545 2220010111:753 0211021122:246
1210210002:5028 1201020110:7506 1012011110:6099 1102211121:408
0011010102:12240 2120011011:1092 2010010001:4521 1201111112:4167
2002221222:408 0022221202:3437 2102221011:1830 1122111202:16 1100201122:1461
'''),
(39, '1/20', 131434191810,
 [1, -12, -12, 79, 79, -103, -103, 53, 53, -12, -12, 1, 1],
 '''
102112020210:1219154721 210001220021:628594164 010210102001:3382052226
002202222201:3675239955 010111122122:1728477428 021100020010:2940379970
100111000001:3965588986 200211122010:2906351488 010011102102:4734733751
111222202210:3564520628 020102000120:5735642133 100112021101:4034972285
221111021112:2788054669 020022212001:7937340702 202022002101:637385244
111122121020:8713871385 221201220112:3725747503 002210102000:5800572994
020110100222:1417269602 211222220111:3281813885 102222220002:9389547438
010021111110:1490195401 202011220110:10254004493 110100110021:725348383
222201122022:6875879692 112000100021:1660697195 210002210000:5226070479
000100102220:4969493108 110010001010:8367080149 221110110102:3745625330
212211012222:4187656967 122021112212:1724829456
'''),
]

def trim(a):
 while len(a)>1 and a[-1]==0:a.pop()
 return a

def divide(a,b):
 a=a[:];b=trim(b[:]);assert b[-1]==1
 quo=[0]*max(1,len(a)-len(b)+1)
 while len(a)>=len(b) and any(a):
  j=len(a)-len(b);v=a[-1];quo[j]=v
  for i,x in enumerate(b):a[i+j]-=v*x
  trim(a)
 return trim(quo),trim(a)

@lru_cache(None)
def cyclotomic(n):
 a=[-1]+[0]*(n-1)+[1]
 for d in range(1,n):
  if n%d==0:
   a,r=divide(a,list(cyclotomic(d)));assert r==[0]
 return tuple(a)

def verify(q,c_text,den,P,pairs):
 pairs=[s.split(':') for s in pairs.split()]
 seeds=[list(map(int,z)) for z,a in pairs]
 nums=[int(a) for z,a in pairs]
 assert q%2 and q>=3
 D=len(P)-1;assert P[-1]==1
 # Verify the real polynomial by its exact cyclotomic lift.
 lift=[0]*(2*D+1)
 for a,coef in enumerate(P):
  for b in range(a+1):lift[D+2*b-a]+=coef*math.comb(a,b)
 assert tuple(trim(lift))==cyclotomic(2*q)
 # The lift implies P(alpha)=0 and degree(P)=phi(2q)/2.
 f=[[2],[0,1]]
 for j in range(2,q):
  a=[0]+f[-1];b=f[-2]
  a += [0]*max(0,len(b)-len(a))
  for i,v in enumerate(b):a[i]-=v
  f.append(divide(trim(a),P)[1])
 inds=list(range(1,q,2));pos={j:i for i,j in enumerate(inds)};m=len(inds)
 V=[f[j]+[0]*(D-len(f[j])) for j in inds]
 perms=[]
 for a in range(1,q):
  if math.gcd(a,2*q)==1:
   pp=[]
   for j in inds:
    z=(a*j)%(2*q);z=min(z,2*q-z);pp.append(pos[z])
   perms.append(pp)
 assert len(perms)==D
 assert len(seeds)==len(nums) and all(a>0 for a in nums) and sum(nums)==den
 assert all(len(z)==D and all(a in (0,1,2) for a in z) for z in seeds)
 c=Fraction(c_text);assert c>Fraction(1,q-1)
 # Accumulate original (not orbit-reduced) moments.
 diff=[[0]*m for _ in range(m)];summ=[[0]*m for _ in range(m)]
 for seed,num in zip(seeds,nums):
  e=[sum(a*b for a,b in zip(v,seed))%3 for v in V]
  for pp in perms:
   ee=[e[j] for j in pp]
   for i in range(m):
    for j in range(i,m):summ[i][j]+=num*(2 if (ee[i]+ee[j])%3==0 else -1)
    for j in range(i):diff[i][j]+=num*(2 if (ee[i]-ee[j])%3==0 else -1)
 scale=2*D*den
 for i in range(m):
  for j in range(i,m):assert summ[i][j]==0
  for j in range(i):assert Fraction(diff[i][j],scale)==-c
 print(f'q={q}: EXACT PASS; {len(seeds)} seeds x {D} Galois images; c={c} > 1/{q-1}.')


for entry in DATA: verify(*entry)
\end{verbatim}
}


\section{Twice an odd prime power}
\label{sec:twice-prime-powers}

The zero eigenvalue creates an additional first-moment constraint when
$q$ is even. The following construction handles that constraint as well
as the relations among the nonzero eigenvalues.

\begin{theorem}\label{thm:twice-prime-powers}
Let $p\ge5$ be prime and $k\ge2$. There is a bounded, compactly
supported probability density on $[0,\infty)$ for which the averaged
mixing matrix of $P_{2p^k-1}$ is exactly $J/(2p^k-1)$.
\end{theorem}

\begin{proof}
Set
\[
 s=p^k,\quad q=2s,\quad h=(p-1)/2,\quad
 r=(p^{k-1}-1)/2,\quad \ell=r+h,\quad c_0=(2s-1)^{-1}.
\]
For $c$ near $c_0$ define
\[
 B=\sqrt{p(1-(p-1)c)},\quad \beta=B/p,\quad
 a=\frac{pc}{1-(p-1)c};\qquad a\beta^2=c.
\tag{EP1}
\]
At $c=c_0$, $a=(4r+1)^{-1}$. We shall construct the required
coherences for $c>c_0$ sufficiently close, then use a positive Haar
component and the exact finite-time transfer proved earlier.

\paragraph{The two independent frequency sectors.}
Write $\zeta=e^{\pi i/(2s)}$ and
$\lambda_j=\zeta^j+\zeta^{-j}$, $1\le j<2s$.
The field automorphism $\zeta\mapsto-\zeta$ fixes every even-index
frequency and negates every odd-index frequency. It is indeed an
automorphism, since $\gcd(2s+1,4s)=1$. Thus the rational spans of
the two sectors intersect only at zero. The odd sector contains the
zero mode $j=s$, whose phase is fixed at one.

After selecting one signed representative from each nonzero pair
$\{\lambda,-\lambda\}$, each sector has $pr+h$ representatives,
partitioned into $r$ disjoint $p$-tuples with signed sum zero and
$h$ free representatives. Here is a direct construction for both
sectors. Put $t=p^{k-1}$. In the even sector, write
$\lambda_{2u}=2\cos(\pi u/s)$ and choose the odd representatives
$1\le u\le s-2$. For each odd $1\le j<t$, take the $p$ exponents
$j+2tv$, $0\le v<p$, modulo $2s$, and reflect into $[0,s]$.
They are distinct odd indices with cosine sum zero. Their residue
classes modulo $2t$, up to sign, show that the blocks are disjoint;
the unused representatives are $t,3t,\ldots,(p-2)t$.
These representatives span the real field
$\mathbb Q(2\cos(\pi/s))$, of degree $\varphi(s)/2=ht$.
Indeed their sum is $1$, and every cosine at an integral multiple
of $\pi/s$ is their signed representative or the constant $2$;
powers of $2\cos(\pi/s)$ expand into such cosines.

In the odd sector, for each odd $1\le j<t$, sum the
$p$ cosines with exponents $j+4tu$, $0\le u<p$, modulo $4s$.
Their sum is zero. Reflect the exponents into $[0,2s]$ and then
into $[0,s]$, recording a minus sign at the second reflection.
This gives $p$ distinct odd representatives. Residue classes modulo
$2t$, up to sign, prove that the resulting tuples are disjoint.
The unused representatives are $t,3t,\ldots,(p-2)t$.

The odd frequency span is the negative eigenspace of the displayed
field involution. Odd powers of $2\cos(\pi/(2s))$ span that
eigenspace and expand into odd cosines; consequently its dimension
is $\varphi(s)/2=ht$. The even sector has the same dimension.
Since $pr+h-ht=r$, the displayed relations exhaust the rational
relations in each sector. Their integer lattices are saturated:
the rows have disjoint supports and coefficients $\pm1$.
Conjugating representative coordinates where necessary makes every
block relation a product equal to one. There are no cross-sector
relations. The time-flow closure is therefore the direct product
of two tori, each consisting of $r$ product-one $p$-blocks and
$h$ free phases. The character criterion also gives equidistribution
of the actual nonnegative-time flow on this torus.

\paragraph{An exchangeable root-label construction.}
For either $\mu=0$ or $\mu=-a\beta$, we construct $\ell$ exchangeable
$p$th roots $Y_i$ satisfying
\[
 \mathbb E Y_i=\mu,\qquad \mathbb E Y_i^2=-a,\qquad
 \mathbb E Y_i\overline{Y_j}=\mathbb E Y_iY_j=-a\quad(i\ne j).
\tag{EP2}
\]
Here and below all prescribed moments are real. We first treat $k\ge3$. Since
$\ell\equiv-1\pmod p$, take $(\ell+1)/p$ complete sets of the
$p$th roots, remove one copy of $1$, and uniformly permute the
remaining multiset. Its sum and its sum of squares both equal $-1$.
Rotate it by a common $p$th root $Z$, then independently conjugate
each coordinate with probability $1/2$. Put $z=a\ell$ and
$b=-\ell\mu$, and assign to $Z=e^{i\theta}$ the mass
\[
 \frac{1+2b\cos\theta+2z\cos(2\theta)}p.
\tag{EP3}
\]
Fourier orthogonality on the $p$th roots, using $p\ge5$, gives
$\mathbb EZ=b$ and $\mathbb EZ^2=z$. The constructed vector has
first moment $-b/\ell=\mu$, second moment $-z/\ell=-a$, and
both off-diagonal moments equal $-d$, where
\[
 d=\frac1{2\ell}-\frac{a}{\ell-1}.
\tag{EP4}
\]
Indeed, before the independent conjugations the sum of real parts
is $-\cos\theta$. Its expected square is $(1+z)/2$, whereas the
expected sum of squared real parts is $(\ell-z)/2$. Subtracting
and dividing by $\ell(\ell-1)$ proves (EP4).

These are probability masses with a strict margin at $c_0$.
Since $r\ge2h(h+1)$ and $h\ge2$,
\[
 z=\frac{r+h}{4r+1}\le\frac27,\qquad
 \beta^2<\frac1p\le\frac15.
\]
For $x=\cos\theta$, the numerator in (EP3) is at least
\[
 1-2z-\frac{b^2}{4z}\ge\frac{29}{70}>0.
\]
An iid root vector with single-coordinate masses
\[
 \frac{1+2\mu\cos\theta-2a\cos(2\theta)}p
\tag{EP5}
\]
has the same first and second moments and off-diagonal moments
$\mu^2$. Its masses are positive: at $c_0$, $a\le1/49$ and
$|\mu|\le a/\sqrt5$. Finally $d>a$, because this is equivalent to
\[
 (4r+1)(\ell-1)>2\ell^2.
\]
The difference is $2r^2-3r+h-1-2h^2$; it increases with $r$
and, at $r=2h(h+1)$, equals $8h^4+16h^3-5h-1>0$.
Mix the dependent law with the iid law, giving the former weight
$(a+\mu^2)/(d+\mu^2)\in(0,1)$. This proves (EP2).
All inequalities are strict, so the construction persists for
$c>c_0$ sufficiently close.

\paragraph{The square case.}
It remains to construct (EP2) when $k=2$, for which $r=h$ and
$\ell=p-1$. At $c_0$,
\[
 a=\frac1{2p-1},\qquad \beta^2=\frac{2p-1}{2p^2-1}.
\]
Choose the integer $j$ nearest to $p/6$ and put
\[
 u=\cos(2\pi j/p),\quad \sigma=1-2u,\quad
 \tau=3-4u^2,\quad D=1+(\tau-\sigma^2)/2.
\]
Since $p\equiv\pm1\pmod6$, the angle lies in
$[2\pi/7,2\pi/5]$. Hence $3/10<u<5/8$ and
\[
 \sigma^2<4/25,\qquad \tau>23/16,\qquad 27/16<D<9/4.
\tag{EP7}
\]
The upper bound for $u$ follows, for example, from the polynomial
$8u^3+4u^2-4u-1$ for $\cos(2\pi/7)$, which is increasing
between $1/2$ and $5/8$ and positive at $5/8$.

Use two root multisets of length $\ell$. Let $M_0$ be all $p$th
roots with $1$ removed. Let $M_1$ be all $p$th roots with
$1,e^{2\pi i j/p},e^{-2\pi i j/p}$ removed and two copies of $1$
added. Their sums and square sums are respectively $(-1,-1)$ and
$(\sigma,\tau)$. Uniformly permute the chosen multiset, rotate it
by a common root, then independently conjugate each coordinate.
Define
\[
 V=\ell/2-a\ell(\ell-1/2),\quad
 w=\frac{1-\ell+2a\ell^2}{D},\quad
 z_0=\frac{a\ell-\tau w/2}{1-w},\quad
 b_0=\frac{-\ell\mu}{1-w}.
\]
Choose $M_1$ with probability $w$, using rotation masses
$(1-\cos(2\theta))/p$. Choose $M_0$ with probability $1-w$,
using rotation masses
\[
 \frac{1+2b_0\cos\theta+2z_0\cos(2\theta)}p.
\tag{EP8}
\]
The first two total moments are
\[
 \mathbb E\sum_iY_i=-(1-w)b_0=\ell\mu,\qquad
 \mathbb E\sum_iY_i^2=-(1-w)z_0-w\tau/2=-a\ell.
\]
Moreover, writing $C_i=\operatorname{Re}Y_i$, we have
\[
 \mathbb E\Big(\sum_i C_i\Big)^2
 =\frac{(1-w)(1+z_0)}2+\frac{w\sigma^2}{4}
 =\frac{1+a\ell-wD}{2}=V,
 \qquad \mathbb E\sum_i C_i^2=\frac{\ell(1-a)}2.
\]
Exchangeability and the independent conjugations make both
 off-diagonal moments equal to
$[V-\ell(1-a)/2]/[\ell(\ell-1)]=-a$. Thus (EP2) holds once
these are verified to be probability laws.

At $c_0$ put
\[
 H=(2p-1)D-p,\qquad N=(p-1)D-p\tau/2.
\]
Then $w=p/[(2p-1)D]\in(0,1)$, $z_0=N/H$, and
$|b_0|\le\beta(p-1)D/H$. Also
\[
 N=\frac{p(2-\sigma^2)}2-D>\frac{23p}{25}-\frac94>0,
 \qquad 1-2z_0=\frac{D+p(\tau-1)}H>0.
\]
The numerator of (EP8) is bounded below on the whole circle by
$1-2z_0-b_0^2/(4z_0)$. This is positive provided
\[
 4N[D+p(\tau-1)]>\beta^2(p-1)^2D^2.
\]
Since $\beta^2<1/p$, (EP7) reduces this to the sufficient inequality
\[
 4p\Big(\frac{23p}{25}-\frac94\Big)(7p+27)>81(p-1)^2.
\]
The left side minus the right side is
$644p^3/25-1116p^2/25-81p-81$, which equals $1618$ at $p=5$
and has positive increasing derivative for $p\ge5$. Therefore
(EP8) has a strict positive margin. The other rotation law is
nonnegative and fixed. All strict conditions persist for
$c>c_0$ sufficiently close, completing (EP2) also for $k=2$.

\paragraph{Templates with two power sums.}
There exist $p$ unit phases $v_1,\ldots,v_p$ such that
\[
 \prod v_i=1,\qquad \sum v_i=B,\qquad \sum v_i^2=B^2/p.
\tag{EP6}
\]
Use a template $(1,e^{i\theta_1},e^{-i\theta_1},\ldots,
 e^{i\theta_h},e^{-i\theta_h})$. With $x_j=\cos\theta_j$,
it suffices to prescribe
\[
 \sum x_j=(B-1)/2=:S,\qquad
 \sum x_j^2=h(1-c)/2=:Q.
\]
For the chosen $c$ we have $0<c<1/p$ and $1<B<\sqrt p$.
On the connected cube slice $\sum x_j=S$, the minimum square sum
is $S^2/h<1/2$, while $Q>h(1-1/p)/2\ge4/5$.
An extreme point of this slice has at least $h-1$ coordinates
in $\{-1,1\}$, so its square sum is at least $h-1\ge h/2>Q$.
Continuity supplies the desired template.

\paragraph{The phase law and the zero mode.}
Use the first $r$ labels in (EP2) for the $p$-blocks. In block $a$,
uniformly permute the template (EP6) multiplied by $Y_a$.
The product is $Y_a^p=1$, the conditional coordinate mean is
$\beta Y_a$, and both the diagonal and distinct holomorphic
second moments are $\beta^2\mathbb EY_a^2=-c$.
The distinct internal Hermitian moment is
$(B^2-p)/(p(p-1))=-c$.

For each of the $h$ free phases, use one remaining label: retain
it with probability $\beta$, and otherwise use an independent
$p$th root of mean zero and second moment $a\beta$.
The latter is supplied by masses
$(1+2a\beta\cos(2\theta))/p$, which are positive.
Its resulting first moment is $\beta\mu$, and its second moment is
\[
 \beta(-a)+(1-\beta)a\beta=-a\beta^2=-c.
\]
All distinct cross-block, block/free and free/free Hermitian and
holomorphic moments are also $\beta^2(-a)=-c$.
Thus every representative has first moment $\beta\mu$.

Choose $\mu=0$ for the even sector and $\mu=-a\beta$ for the odd
sector, and take the two sector laws independently. The first
moments are respectively zero and $-c$. Every same-parity
coherence between distinct modes is $-c$, including a mode and
its negative partner, and including the fixed zero mode in the odd
sector. Every opposite-parity coherence is zero by independence.
Independent conjugations used in orienting the signed relations
preserve all these real target moments.

We have constructed the full forced-coherence target at $c>c_0$
on the exact time-flow torus. Mixing with Haar measure in proportion
$c_0/c$ gives the target at $c_0$ and a positive Haar component.
All nonzero frequency-difference characters have Haar mean zero.
The previously proved smoothing, finite-character correction and
finite-time Gram correction yield an actual bounded compactly
supported density on $[0,\infty)$. The path spectral identity now
gives $J/(q-1)$.
\end{proof}

Together with the known positive result for twice a prime and the
known parity-collision obstruction, this classifies all moduli
$q=2p^k$ with $p$ odd prime and $k\ge1$: the negative cases are
exactly $p=3$, $k\ge2$. The next section extends the positive
construction to higher powers of two when $k\ge2$. General
mixed-prime moduli remain outside this theorem. This is a
classification of one infinite family, not an all-path classification.


\section{Higher powers of two times an odd prime power}
\label{sec:higher-two-powers}

The same templates also accommodate a larger power of two in the
path modulus. The new ingredient is a root-label law for the changed
numbers of relation blocks and free phases.

\begin{theorem}\label{thm:higher-two-powers}
Let $p\ge5$ be prime, $k\ge2$ and $A\ge2$. There is a bounded,
compactly supported probability density on $[0,\infty)$ whose
averaged mixing matrix on $P_{2^A p^k-1}$ is exactly
$J/(2^A p^k-1)$.
\end{theorem}

\begin{proof}
Put $R=2^{A-2}p^{k-1}$, so $q=4pR$, and use $B,\beta,a$ as in
(EP1), now with $c_0=(4pR-1)^{-1}$. At $c_0$,
\[
 a=(4R-1)^{-1},\qquad \beta^2=(4R-1)/(4pR-1)<1/p.
\]
In particular $p\mid R$ and $R\ge p$.

\paragraph{The relation lattice.}
Let $\zeta=e^{2\pi i/(8pR)}$. The involution
$\zeta\mapsto-\zeta$ separates even and odd index frequency spans.
Both have dimension $(p-1)R$. Indeed, the even span is the real
field of the primitive $q$th root, of degree $\varphi(q)/2$;
the constant belongs to this span by the sum of the nontrivial
$p$th roots, and cosine expansions of powers give the whole field.
The odd span is the negative eigenspace, since odd powers of
$\zeta+\zeta^{-1}$ expand into odd-index cosines.

For each $1\le j<2R$, sum the $p$ cosines with exponents
$j+8Rv$, $0\le v<p$, modulo $8pR$. Reflect into $[0,4pR]$
and then $[0,2pR]$, recording the sign of the second reflection.
These are disjoint $p$-tuples: their residues modulo $4R$, up to
sign, distinguish them, and $0<j<2R$ excludes a repeated or zero
representative. Odd $j$ give $R$ blocks covering all $pR$ odd
representatives. Even $j$ give $R-1$ blocks, leaving the $p-1$
free representatives $2R,4R,\ldots,2(p-1)R$. The zero mode is
$j=2pR$, which is even. The dimension counts give exactly $R$
and $R-1$ relations respectively. Disjoint support and coefficients
$\pm1$ make the displayed relation lattices saturated. Thus the
exact time-flow torus consists of these product-one blocks and
free phases, with independent parity sectors.

\paragraph{Odd-sector labels.}
Take $R/p$ complete sets of $p$th roots, uniformly permute them,
and independently conjugate each coordinate. The first and diagonal
holomorphic moments vanish, while both off-diagonal moments are
$-1/[2(R-1)]$. Mix this law with weight $w=2a(R-1)$ and an iid
law with masses
\[
 \frac{1-2a\cos(2\theta)/(1-w)}p.
\]
The result has zero first moments and every required second moment
$-a$. At $c_0$, $0<w<1$ and $a/(1-w)=1/(2R+1)<1/2$,
so this remains valid for $c>c_0$ sufficiently close.

\paragraph{Even-sector labels away from two exceptional cases.}
Set $\ell=R+p-2$ and $\mu=-a\beta$. Take $(\ell+2)/p$ full
root sets, remove two copies of $1$, uniformly permute, rotate
by a common root of first moment $b=-\ell\mu/2$ and second
moment $z=a\ell/2$, then independently conjugate the coordinates.
The rotation masses are
\[
 \frac{1+2b\cos\theta+2z\cos(2\theta)}p.
\tag{HP1}
\]
The multiset sum and square sum are both $-2$. Thus first moments
are $\mu$, diagonal holomorphic moments are $-a$, and both
off-diagonal moments are $-d$, with
\[
 d=\frac{(1-3a)\ell-4}{2\ell(\ell-1)}.
\]
For the last formula, the squared total real part has expectation
$2(1+z)$ and the sum of squared real parts has expectation
$(\ell-2z)/2$. Mix this law with the iid law (EP5), using dependent
weight $(a+\mu^2)/(d+\mu^2)$. This gives every second moment
$-a$ whenever $d>a$.

At $c_0$, this inequality is equivalent to
\[
 F:=R^2-9R-p^2+3p>0.
\tag{HP2}
\]
For $R\ge2p$, $F$ increases with $R$ and at $R=2p$ equals
$3p(p-5)$. Hence (HP2) holds except at $(p,R)=(5,10)$.
The allowed $R$ values have no value strictly between $p$ and
$2p$. For $R\ge2p$, $z<3/16$ and $|b|=z\beta$, so the numerator
of (HP1) is greater than $197/320$. Also $a\le1/39$ and
$|\mu|<a/\sqrt5$, so (EP5) is strictly positive. All these
conditions persist for $c>c_0$ close.

\paragraph{The cases $R=p$ and $(p,R)=(5,10)$.}
Choose the integer $j$ nearest $p/4$ and put
\[
 u=\cos(2\pi j/p),\quad v=u^2=\sin^2(\pi/(2p)),\quad
 T=2-4v,\quad D=3(1-v).
\]
Since $\pi^2<10$, $0<v<5/(2p^2)$. Let $M_0$ be the preceding
multiset with sum and square sum $-2$. From $M_0$, remove one
copy of each of $e^{2\pi i j/p},e^{-2\pi i j/p}$ and add two
copies of $1$, producing $M_1$ with sum $-2u$ and square sum $T$.
Both multisets have $\ell$ entries. Define
\[
 V=\ell/2-a\ell(\ell-1/2),\quad
 w=\frac{2+a\ell-V}{D},\quad
 z_0=\frac{a\ell-wT/2}{2(1-w)},\quad
 b_0=\frac{-\ell\mu}{2(1-w)}.
\]
Use $M_1$ with probability $w$, rotating with masses
$(1-\cos(2\theta))/p$. Use $M_0$ with probability $1-w$,
rotating with masses
\[
 \frac{1+2b_0\cos\theta+2z_0\cos(2\theta)}p.
\tag{HP3}
\]
Uniformly permute and independently conjugate in each case.
The total first moment is $-2(1-w)b_0=\ell\mu$, the total
holomorphic square moment is
$-2(1-w)z_0-wT/2=-a\ell$, and the squared total real part has
expectation
\[
 2(1-w)(1+z_0)+wv=2+a\ell-wD=V.
\]
Consequently both off-diagonal moments equal $-a$.

For $R=p$, let $H=(2p-1)-(4p-1)v$. At $c_0$ the parameters are
\[
 w=\frac{2p}{(4p-1)(1-v)}\in(0,1),\qquad
 z_0=\frac{(p+1)v-1}{H}<0,\qquad
 |b_0|=\frac{(p-1)(1-v)\beta}{H}.
\]
Here $H>0$ and $z_0<0$ follow from $v<5/(2p^2)$ and $p\ge5$.
The numerator of (HP3) is at least
\[
 1+2z_0-2|b_0|
 =\frac{(1-v)[2p-3-2(p-1)\beta]}H>0,
\]
since $\beta<1/\sqrt p<1/2$. All necessary inequalities are
strict and persist for $c>c_0$ close.
For $(p,R)=(5,10)$, (HP2) is an equality and $w(c_0)=0$.
At $c_0$, (HP3) is the strictly positive law (HP1). Since
$dw/da=\ell(\ell+1/2)/D>0$, the same construction is valid
with $0<w<1$ for $c>c_0$ sufficiently close.

\paragraph{Assembly and actual time.}
Apply template (EP6) to the blocks. For each even free phase,
retain its label with probability $\beta$ and otherwise use
independent root noise of first moment zero and second moment
$a\beta$, exactly as in the previous theorem. All distinct
same-parity coherences are $-c$, including the negative partners;
first moments are zero in the odd sector and $-c$ in the even
sector, giving the correct coherences with the fixed even zero
mode. Independent sectors give zero opposite-parity coherences.
Thus the full target holds at some $c>c_0$. The positive Haar
mixture and exact finite-time transfer yield the claimed density.
\end{proof}

This extension leaves moduli $2^A p$ with $A\ge2$ and general
mixed odd-prime factors outside the construction. It is not an
all-path classification.


\section{What the constructions do and do not resolve}

The local-potential problem is a finite auxiliary problem with an exact
answer: the displayed incidence rules admit the rational certificate
with constant $875/68$. The profile constraints sharpen its consequence
from twelve vertices to eleven. Applied to the original radius
inequality, this proves the complete degree-seven case for graphs of
arbitrary order. It does not remove the degree restriction. In
particular, a proof of the auxiliary certificate problem is not being
presented as a resolution of the unrestricted Fajtlowicz conjecture.

The phase-moment problem has two additional mechanisms: a quadratic
torus density couples a dependent phase to the remaining modes, and
exchangeable root multisets satisfy the first moments imposed by a
fixed zero mode. Strict feasibility allows an exact correction along
the actual time flow. These constructions settle the stated path
families with the observation rule unchanged. A complete classification
of all path moduli remains outside the results.

The source search checked the repository snapshot, the radius papers
\cite{CPS,Doig2023}, the quantum-walk problem survey
\cite{CoutinhoGuo2024}, and Kang's arithmetic classification
\cite{Kang2026}, together with targeted searches for higher-degree
radius bounds and the path families treated here. No identical solution
was identified in the inspected material. This is a bounded search,
not a proof that no earlier equivalent argument exists. Neither a
historically unprecedented principle nor a ranking above all previous
corpus results follows from it. The manuscript is a research draft
supported by internal independent checks; it has not undergone external
peer review or proof-assistant verification.

\appendix
\section{Reproducibility from this single source}
\label{app:repro}

The TeX source compiles without external figures, bibliography files,
or data files. It also contains, after the document terminator, a
compressed text bundle with the following contents:
\begin{itemize}[nosep]
\item the full rational graph certificate and its Python verifier;
\item the C finite-core checker and an independent Python checker;
\item exact quantum certificates for $q=33,39$ and their verifier;
\item the recorded finite graph checks and generator provenance.
\end{itemize}
Compression is only a transport format: all mathematical data are
ordinary decimal integers or rational strings. Python's standard
library suffices to extract and verify the rational certificates.

Save the following as \texttt{extract.py}, then run
\texttt{python3 extract.py paper.tex}, replacing \texttt{paper.tex}
with the actual source filename. The script creates a fresh
\texttt{repro} directory beside the source.
\begin{lstlisting}[language=Python]
import base64, json, pathlib, sys, zlib
source = pathlib.Path(sys.argv[1])
s = source.read_text(encoding="utf-8")
begin = "%" + "BEGIN_EMBEDDED_REPRO_BUNDLE"
end = "%" + "END_EMBEDDED_REPRO_BUNDLE"
payload = s.split(begin, 1)[1].split(end, 1)[0]
encoded = "".join(line[1:].strip() for line in
                  payload.splitlines() if line.startswith("%"))
files = json.loads(zlib.decompress(base64.b64decode(encoded)))
out = source.parent / "repro"
out.mkdir(exist_ok=False)
for name, content in files.items():
    assert pathlib.Path(name).name == name
    (out / name).write_text(content, encoding="utf-8")
print("Extracted", len(files), "files to", out)
\end{lstlisting}

From the extracted directory, the two fast exact checks are:
\begin{lstlisting}[language=bash]
python3 graph_verify_compact.py graph_degree7_compact_certificate.json
python3 quantum_compute_exceptional_verify.py
\end{lstlisting}
The first must report $45$ types, $802$ allowed pairs, $896$ profiles,
$895$ profiles covered by the fourteen ordinary dual vectors, and
residual lower bound $103/84$. The second checks every required
Hermitian and holomorphic moment for both exceptional quantum cases.
Neither verification needs a linear optimization package.

To reproduce the exhaustive graph step, install the cited version of
\texttt{geng} from nauty \cite{Nauty}. Compile the bundled C checker
and run the following commands from the extracted directory:
\begin{lstlisting}[language=bash]
cc -O3 -std=c11 graph_corechk.c -o graph_corechk
for n in 4 5 6 7 8 9; do
    geng -C -d2 -q "$n" | ./graph_corechk
done
geng -C -d2 -D7 -q 10 | ./graph_corechk
geng -C -d2 -D7 -q 11 | ./graph_corechk
python3 graph_corecheck_python.py /absolute/path/to/geng
\end{lstlisting}
The C program's pass/fail decision uses only integers; its decimal
minimum-slack display is diagnostic. The independent Python program
uses exact fractions and checks orders four through nine. The two
higher-order minimum-slack witnesses were also checked with exact
fractions. The generator supplies completeness up to graph isomorphism;
the checker does not itself certify the generator. The C parser is
intended for these small, header-free graph6 records only.

The $q=33,39$ verifier is printed in full in
Section~\ref{sec:exact-33-39}. The remaining quantum proofs are analytic
and require no finite search. Their use of continuity is justified by
explicit strict inequalities at the target parameter; in the single
boundary case $q=200$, the mixture weight increases from zero on the
required side while the other probability masses remain positive.

\section{The global graph potential}
For direct inspection, the first graph certificate is printed below.
Each row gives $t=(d,p,l)$ and the multipliers in
\eqref{eq:g7-vcert}--\eqref{eq:g7-ecert2}, with $C=875/68$.
The profile certificates and residual vector are in the embedded
bundle. Their exact coefficient inequalities are specified in
Lemma~\ref{lem:g7-profiles} and checked by the bundled program.
\begin{center}\small
\renewcommand{\arraystretch}{1.2}
\begin{tabular}{crr@{\qquad}crr@{\qquad}crr}
\toprule $t$&$\alpha_t$&$\beta_t$&$t$&$\alpha_t$&$\beta_t$&$t$&$\alpha_t$&$\beta_t$\\\midrule
$(3,1,1)$ & $0$ & $-\frac{671}{408}$ & $(5,1,2)$ & $\frac{61}{136}$ & $-\frac{125}{272}$ & $(6,1,6)$ & $\frac{1}{5}$ & $-\frac{137}{170}$\\[2pt]
$(4,0,1)$ & $\frac{875}{272}$ & $-\frac{201}{272}$ & $(5,1,3)$ & $\frac{313}{400}$ & $-\frac{787}{850}$ & $(7,0,1)$ & $\frac{137}{170}$ & $\frac{11}{340}$\\[2pt]
$(4,0,2)$ & $\frac{919}{544}$ & $-\frac{647}{544}$ & $(5,1,4)$ & $\frac{201}{272}$ & $-\frac{253}{136}$ & $(7,0,2)$ & $\frac{57}{136}$ & $\frac{11}{340}$\\[2pt]
$(4,0,3)$ & $\frac{671}{408}$ & $-\frac{535}{136}$ & $(5,1,5)$ & $\frac{1}{4}$ & $-\frac{253}{136}$ & $(7,0,3)$ & $\frac{74}{255}$ & $\frac{11}{340}$\\[2pt]
$(4,0,4)$ & $\frac{1}{4}$ & $-\frac{2743}{816}$ & $(6,0,1)$ & $\frac{253}{136}$ & $-\frac{117}{680}$ & $(7,0,4)$ & $\frac{373}{1360}$ & $-\frac{11}{340}$\\[2pt]
$(4,1,1)$ & $0$ & $-\frac{201}{272}$ & $(6,0,2)$ & $\frac{787}{850}$ & $-\frac{181}{850}$ & $(7,0,5)$ & $\frac{181}{850}$ & $-\frac{11}{340}$\\[2pt]
$(4,1,2)$ & $\frac{1007}{816}$ & $-\frac{647}{544}$ & $(6,0,3)$ & $\frac{53}{85}$ & $-\frac{74}{255}$ & $(7,0,6)$ & $\frac{117}{680}$ & $-\frac{11}{340}$\\[2pt]
$(4,1,3)$ & $\frac{671}{408}$ & $-\frac{2743}{816}$ & $(6,0,4)$ & $\frac{125}{272}$ & $-\frac{57}{136}$ & $(7,0,7)$ & $\frac{1}{7}$ & $-\frac{11}{340}$\\[2pt]
$(4,1,4)$ & $\frac{1}{3}$ & $-\frac{2743}{816}$ & $(6,0,5)$ & $\frac{107}{272}$ & $-\frac{263}{272}$ & $(7,1,1)$ & $0$ & $-\frac{19}{136}$\\[2pt]
$(5,0,1)$ & $\frac{175}{68}$ & $-\frac{107}{272}$ & $(6,0,6)$ & $\frac{1}{6}$ & $-\frac{137}{170}$ & $(7,1,2)$ & $\frac{1}{6}$ & $-\frac{19}{136}$\\[2pt]
$(5,0,2)$ & $\frac{647}{544}$ & $-\frac{125}{272}$ & $(6,1,1)$ & $0$ & $-\frac{467}{2040}$ & $(7,1,3)$ & $\frac{79}{408}$ & $-\frac{125}{816}$\\[2pt]
$(5,0,3)$ & $\frac{404}{425}$ & $-\frac{787}{850}$ & $(6,1,2)$ & $\frac{787}{850}$ & $-\frac{2093}{5100}$ & $(7,1,4)$ & $\frac{49}{272}$ & $-\frac{125}{816}$\\[2pt]
$(5,0,4)$ & $\frac{389}{544}$ & $-\frac{253}{136}$ & $(6,1,3)$ & $\frac{787}{850}$ & $-\frac{7271}{10200}$ & $(7,1,5)$ & $\frac{49}{272}$ & $-\frac{1}{6}$\\[2pt]
$(5,0,5)$ & $\frac{1}{5}$ & $-\frac{253}{136}$ & $(6,1,4)$ & $\frac{125}{272}$ & $-\frac{1009}{1632}$ & $(7,1,6)$ & $\frac{117}{680}$ & $-\frac{1}{6}$\\[2pt]
$(5,1,1)$ & $\frac{2743}{816}$ & $-\frac{107}{272}$ & $(6,1,5)$ & $\frac{107}{272}$ & $-\frac{409}{408}$ & $(7,1,7)$ & $\frac{1}{6}$ & $-\frac{1}{6}$\\[2pt]
\bottomrule\end{tabular}\end{center}

\begingroup\small\setlength{\parskip}{0pt}
\begin{thebibliography}{9}
\setlength{\itemsep}{3pt}
\bibitem{CPS}
M.~Cygan, M.~Pilipczuk, and R.~\v{S}krekovski,
\emph{On the inequality between radius and Randi\'c index for graphs},
MATCH Communications in Mathematical and in Computer Chemistry
\textbf{67} (2012), 451--466.
\url{https://match.pmf.kg.ac.rs/electronic_versions/Match67/n2/match67n2_451-466.pdf}.

\bibitem{Doig2023}
M.~I.~Doig, \emph{Block eccentricity, a radius bound, and an application
to the Randi\'c index}, arXiv:2309.11613v1, 20 September 2023.
\url{https://arxiv.org/abs/2309.11613}.

\bibitem{Corpus}
QEEDDD, \emph{Verifiable Mathematical Reasoning Corpus},
102 numbered research drafts, snapshot
\texttt{ecd87de3348dcb706a76bb12fde722d7a348a012}, 2 November 2026.
\url{https://github.com/QEEDDD/Verifiable-Mathematical-Reasoning-Corpus}.

\bibitem{Corpus087}
VMRC, item [087], \emph{Spanning trees of small diameter and the
Randi\'c index}, file
\texttt{[087] fajtlowicz\_radius\_core\_reduction.tex},
unrefereed research draft in the snapshot of \cite{Corpus}.

\bibitem{CoutinhoGuo2024}
G.~Coutinho and K.~Guo,
\emph{Selected open problems in continuous-time quantum walks},
Special Matrices \textbf{12} (2024), 20240025.
\href{https://doi.org/10.1515/spma-2024-0025}{doi:10.1515/spma-2024-0025}.

\bibitem{Kang2026}
M.~Kang, \emph{Arithmetic of Bohr Frequencies Governs Uniform Mixing
in Randomly Timed Quantum Spin Chains}, arXiv:2609.14463v1,
13 September 2026.
\url{https://arxiv.org/abs/2609.14463v1}.

\bibitem{Nauty}
B.~D.~McKay and A.~Piperno,
\emph{Practical Graph Isomorphism, II}, Journal of Symbolic Computation
\textbf{60} (2014), 94--112.
\href{https://doi.org/10.1016/j.jsc.2013.09.003}{doi:10.1016/j.jsc.2013.09.003}.
Software: \url{https://users.cecs.anu.edu.au/~bdm/nauty/}.
The runs used nauty 2.9.0, source mirror and commit given in
Section~\ref{sec:g7-main}.
\end{thebibliography}
\endgroup
\end{document}

%BEGIN_EMBEDDED_REPRO_BUNDLE
%eNrtfQtz2zi25l/hZqrbok1JBAgCpBy6N93JvTu725l0TXZ6p2RdFyVSiRKZkvVw4qTTv32/A/AB6mE73bN1a2vTbknkAXDe5+AA
%YKTPT96s0uXbq9t8NZveXU0W18t0sukt754Mnlzqv3/oFmfzNney/M0qz7vr/DYvnEm+2syms0m6yZ3tela8cRbF3PR7dbd5uyic
%9SYtsnSVOfPZeJWu7nqXxf9ap2/ygbPUHQLnCO0SbsipquHKoth7t14Ul8XLhbNYbmbXs0/pZkYUF3MgcxYrZzpfAFS86S4Xs2Lj
%pKvZ5u11vplNnNka7OYZmDHyXRbT1eLama5AAjjWzux6uVhtnH8rAU66dn4pOy3TzVsIU3V5hduyZbYBc4vFvB6+XC2y7WRzWZT3
%b9M1DfWIcW99t74sCFlCKDq47aWrN7dDNnKd2dSZ50UNcy+Yk8/XuSbWubqazub51ZXb+wCBror0Ou+cPE5bJ+4lrLFJE7rpQT3Z
%ukMs9FZ5ml1t8o+bjosur5NhJ/CYx9zR2bCTeUtv7k6h0MyZFc4qLd7kHeFFGrQEqOOjJ93Mm3bmZWcYflm8SIZ4p9aZV6KiXnmx
%vc5X4Kvz2h1cFg7+oz7vGgwzj1Tw2q2a6b/cu/Guk9fDd6MGls7niw95lnTmZ9cXSaaVlyVJbjT2S2fuZaedrMtc9+KXTt699vJT
%fDCSs0KBESWWwYteulzmRdYB/XfUp0gMG+fpeg1tOkWSiNCBU2sLvXCTJPL5ZfFjMvylsyRCc7cmqHVGIpNUr8Hzj85ZMuxSx5u6
%0w67fsWXUZlWyAvHxjR8PZyNRgTR6jCQd6MR8MO6+dSZLObb6+LqNp1v886d977SICi9f1pY6tQYoc73ljpX+Wa7KpxuZ95duqd3
%Q376ftTVYpmbMzY6o4viLOsG5TiwmbwYvu8Wo/MsIe6G/ug813bClelT4qWhswrHuzPEZYekd2vmKeBu8ysjw7rzsmL+fZ4vk+HM
%udeTGjG0tC+HxKMxVnf5NFlvrzsvhwMCuhoKsSoogIOR22Ulu+/XyTrfdIiq2xKAIGfD4uy95uS9p/2kxcYLMumMWHu/1lTemWst
%4jKh+BueLBebvNjM0vnJ6PynBB4xPJkg8SBbbk5GtbP9lKApUqEnI3BRAoeb7XKedzbavTaEG4M3d8t8fTIaJcnry+JVUva5Tped
%X7yPxhM/ll1r0tTfhGapTq+TemMS5tNs2Xntvap0XxI2LpGeGW8Yn/10uu/yFwmrot24bomicrUZHGMniCkAnuaDksgrch/4xyvj
%PE8T3/Si8DjYB5+GkSak9KDL4i/Os/mckjDly7VTLDAHzCnT3Tn5x8l8m+WZMzaz1XqxXU1yZ7zN3uQbTAvVoDp5YaKqM1NgZ6xC
%WEkv6BaBna90h3CnQ7cQdp+6n9zv1y3C3a7038ukUwReIbwi9ArpMV737hbS3e8/na3Wm6Sg9F4F0Md26Lwkn/14YKiOIj1+9FR/
%nAWDSjlVqnzZ+CblxKqZUmMsL4t/JsP5bL0pnfHOOOMd0S9DwfS/yraVS1rI/gksTOgwQoruEOgOoDh0aiz/BPlxmtWmMsBSbWio
%2PxchmyjZsL2o87Aeynz4kdkxS/APFmgmsD04p+v8ut0VqCeqCm91BFVilsSBCaTOvYSmVu7e1rcdeCNHerrfE9cjWdZlhcmKVE+
%+qXz0T39pPn96H1yypB86enMizxl4mzHg+68GlM14p8eFIAptJTirBqlo6kW6IAhK7FhQjPbNdKjOPC9wPyZvP0X58VHiIs+6xlZ
%0a4LzxFrm1Xq6Oy2AorNGszSbOvoommN2mg+m8w2VBq+TGoyNIVoXe7p8Xxd5tCK3EnjMmtKo7rjCVIhjb8s7hLL/dCBhlzlN0iz
%7vmnQ23bMbVdFs+2sPT5eFvbe6f+sbLAyq58dBLq9zkKICt4x4ttkbUnG3RctcsQ3eki8XeC/tm2MtGQZkzM4tomNHNDT0mSNbcc
%KXPl1o5OKhjtxPW4xqbJubb2nmnpk+TZ1vji8GRcQsbbuh9F4i2SrBb9VvuaNT1hRrx1TZa5bWeZTxRqt6NEl5cogyTzBPddt1Tv
%O+99uz8xvzMH7QXqmQmX1eIDJgP39NaYY/HB0KYYeLb1PtGc8KOucxoz6Bi7LWPs9kCMnR3vNdY4a5kNUszVzA+8SKBepymETKUL
%7Wx7vVx3Pp+8mS/G8C8zVw9E6FWQHNMOIKgka1BdDQwQNp2fXA9qOEmz69l6PRsjW1ZpB6Ni6e0k0QET3slilc0KLLnqrtVyIM9O
%qoSgsVZxVHU8Gbz0GiDVxasrLaHhxTiNHrl+m/JQngzKhU3P3FtLivHdJl9jTdF7m3/MZhASCwwzcpNutmD0BL6EDKGzh16/UfZ4
%m09QOi1JudnJF2+GwrDYJJwU/sQrF6uTxSqfvH3fm9Aitej3nZ8AgHnyGzCNbGKwDLTh8nTy1qwnnY7+kB4WptDOObkVvC1zeBcK
%L/LJBjpx8t6bnslOb3Ksabs/uQ4ttZxrmrP0Eusi4RQfmvD6KW5uSbUTVBmLadnF4QPHrGqd1dtF53+4ztPE+Zk+uw5zzpx1PzCC
%z+96zq9O4swn1x3W612nH4HAPZXOegJJkJ2RYGeFrlecp8T2ovf2YgdGyt8FrmgsAcnfqLTeFuvZmyKnLLZxtgE/vyQX0+thAKCM
%8g5NzrNhwEc1hNozJEoC6oZ65O1ilsFSqzVFK4BQe7o6XbufsVBaIzd1ZXDuXOfXNCM+83xvPfuUQ0XP3HPH6r9M1mfsXJMZzzaJ
%7yG+MeNS5MN+HYK/S9j5u6fF+buzM7cGztBp9vTd+QzAz8ifHRqd+O5nGn+6PDvT5Akoz7/QZ7d7Tt3QfHGB2++Z+/kZisnfErZ9
%CjwQ+115M8MA/NWCzhfwBP32ZpJRJVLfp15zPQYbH94ihjq4asCbJP1ufJ4m4/NxsgHici2RnhMBEoQmvQ7GkjYcGD0fIo6g/wbF
%pCC9OLezBWnGyRZbZAHyyPU8YXl8boa+AWAoxSihXQ2tP8PNFGXtukOIKwuYawoCF3Q1SbjMZL0sTDcgoN0Rd0TUoLL/orv4WDjB
%bptZscXMbuxODe55WdyQo7whF0GEEb9UmOHjOWFpW60orYbus1FydTXezubAe7VcLCfIMHAXwF1NW3d5yl2NjTWgBDnBWZ+dNZCL
%5y5o6UutWzOtdjDut9/WF9xmHgrFSNOlUfOv8BSLUQ5Gk+eaUzT92je2/9WbuaezvbE/20JucbOFkNvKXymq0uTZcDs6T8/T75MU
%RcJn6nlrCT/ZfOqkRujbi63r/HwGoiTadnSB99vRD/p6oK/d8y91BaJprn0QXfuguva1bqsGuM75BuANMaMDee2PhptREsNx6nu8
%SAJidJ3nhQ6DNYw3XSV0b+KTKuHKqVZkPlI/DSmmlfh0d5NMV+c35zffJzeQE42/QfS2nDcubFRMv09+N9ibocXUGlqztzc6yc41
%o7/RAOKSPr7UZgezKw4J48NG+ax75LbNWmoilyLSW1LURe6ia3OvLcRP86cr+BSo4LL2t6/zAOc+F3i6tT1WJ4qPNsMfcPMBJD7U
%8tzUTH4YPdWXt3T5QwMdNFBN5Obi+qNLeG9IeVqs649nrJGM7mrhMN/pidVBW5/rKa3/qzWjAXLhOMHpr6cr3SpPf0arPP0V7RzQ
%9W7QzN+uE+reJFR3xT1MmesEY7sYecbbzesy8MoUiPzXMZduB6O6QOj2O7Lnn/5aZ6VpZz1/qnOl+9mkzPWcJrfVZHnXoZzpmSTW
%ZAxguQA2aJUyLvm4ruemncsn//jr3/7ns9d//dtL57u1zpHl4Mr1pmVPJNd8tUIipVJlnXw3n2camS5zynuqKNbzdPI++a4nsDDc
%1DiRnjzq7Wl+PWKSxPmyVwLBGFdmL93s2D9qJ7uErrdj1HwoW9ZmT3ozu87NxvTxrW6zTZflk0VW5n6q0NcJXfao6lh23PMiQfXZ
%oQLApSl4rJdQuhajYgLLCjYYDcY9LA31Xit1nlDPi4uZ+z0zexTNkir0uswrt2zTZEi1hI4o56rpUyCbvNebRfWipFqPFWYNQfBt
%A78tVxa06sLKYJAiPHpplqHhPEV86OutWTq9NwvnatqulGDKO0QrYTIbxYhcVG/NqhIwL3X1mjV/0zMTGyekIAuTdqjUe8rptALT
%06Ak8HJR5Ojys16mYAXCPBSF1BUsevSB3O/uylO4teAkCvBvn966Hu0oEz2K+sYKrXFaD1kyLM7Y6LQ4zzAaeeVmm2/zBNeV7jRq
%DTWKI9gHQ+521GwxZ0gsF8WAPpIMLUgeelC15vygVVounLRc7tNCcwAOq060KKWymSekJOQg3a9eg1NX96goLT3URoY2BjVGuvDY
%mUGs7VDnR6/Jj24tYoVfq9IE7M9ddvZLZ+0FLlayGp+2a7WUpj5YyNOWqr52G//R9+RD4OpKn99cXTlJ4lw+ubqiQvDq6vJJqU+s
%P2gpvMAM3JwPHT0eunxC6xU6zjIb1uvtfGO2MOtNRnsHg/n1/gQK8DoV9F4tYIPOsKbtnXR/OsFbxun95sSjVSDijcrHxXbTGvrX
%Vy88OkdKXq+2ebXvoN0eLjXPUzgheff5JqFc06O3jn3iQVlE76r1DPJ6IyRN7JRzfpvUwVf1oL0FOtmjfV6i0eyhaPoI4ZKBW6NA
%XFJ36mrUR36gwd5thXTZ+5DOkG2qDfllz9iQOEmqrelS08nnk+JkUHgnOt/TEhtUvRNgnV1vr6+01c0KWhNxvZN8ni6xxL1aQ7Ii
%w5AVLa07lma6G/jXlxaZetvV3LrnexsNZYM3nW/Xb21DHNiSOLAOby+/9U4AkQV35mJvQY4Vd7GBhHTquN3kV/nHSb40a/nyXNc+
%SzZbhet8PqUl9wb+jhWptWG41n5wkwSB3oHCRdxz/r5zjKzPmXvV4W05P103R7FH57+qfVtMWke189X2apJCdIrL589eP8OCfHhZ
%dILAc05Yn6kTz4kDGUjavRgyz2HcvLoiKN8CYV5d5pdv1G2ktztOTi4L5vs+XpzzQSSZw/R/3Gd8IGMpHc4Z87nPca+CyKFrRhA2
%CHiM2Yea0Blvg1hw9KdGTigGIubor1s59wdhoAg73YLcgKvYIbI0Fq1RHIjLAsjMPRvIIAwdTUy/DWQocE9jmb5XYeD4mhoNGXAh
%IQunW5/Qhz4HtxiqGfbR3ZfEjUaGe+nHsVPzygbCj0CeutIfHwCn8B1iRcPYgPkQh/D5WmMDEXLmlOi0tEwqUgdpkiQmjFpALQ8f
%BALyE3+kK8IXBb6jpWW6nUliR3OLwUxIKBcGon2pThBrc3Mf5mYBE4FgMYtgy9LqXW1yelOxeXVpy698D8tX02vXA7hmgVQHuYE6
%FArCaU0aa0H9PApjyCgcX+uY9AotBEHE/RAySdIe+nJtbDRIFfIgjmFEYzBtVwimeCSUErCOtp6v8fgDHgs/UHGsfO2STDdpCrEM
%wyiKI/gWcUKIzAAfLsJEFDl+OYC0OxAqwP+BCqE/3VkzhQFBKGEzX2rKhn9yh0GoqIUzRLbGwzVf5BGBiBWH3DArZ4YAWZqrKPJD
%IWVMiIy9iVMVB6DtK5/8hIj6WlHwZEROyIXQDGk9EPVBpFigX+TXOuS08QeB4qEScNjABEjJ6yCMfD9UPI51oBhX1XHLBLSKaCXC
%RmITMgGP4CVBBALMGIbr2IiDKIaJRUAuz4xSdVAwEfsMTT7TEhhVE9znAPpCxIFjyGpXhcg8DEQURAGpiJcjwJGMVBipWFIGYNxk
%GOrPpPRlrEDDeBfX7xANDgTFCRWTG5kg00lDAAeoMsQSq0KPrB/Bv/zIB8M6QzGjDdKdCCWYCuBGOrWU6QfhGSkZQklIQpyXfsS0
%P4oIzhfKOtzqQ3+sGq7LMlpvLTimlEZdow/+hl02wmw7SHvLxbLjHqjEs9ntDLUBHV8P9CIhHQ5GtNlGiMe4rmfwsUFG5fzNdpFg
%haJrQOYZkl36GLtnzN3lJTEthqOiqvqdd0lrIAoTzS5q3sXw3Si5rQrSvSNXcJoOZ2fvRt3k9vQjdSu10IinAUDkenXTZfFf66mq
%Q6WLWz7zcTeZLzaL69mkrIKxTgIbZyRfp8Da6WzIquI/O7A4QkVUfJcl1clT6q2SWqd6MWMRyFAGV+pcJYl5xKNiWT93YDi1lkg3
%3kQ/XoTVS+G98pbpbGVOdvQVlnS99XKOQguFiFnYrHURSI1lCx0S0W5PtraOk1HO0DmMHvHJS+sx1BeKRlcqeNLDHUoJbr7jptC4
%SBBcznNtz1dul53XTxvUHvMXx3oSbpWnc2e5mN8V0Aoux3fObLOuiqhaXShZpvRIgf7U/ob1zHPjYMRW6k0WUFTLN141C5lxY6yU
%DhQJy/A5Vi7jbjo6S2jsKRU+WFtej3UANKKVD4qQ79Aw100Sy4z89MbVMr2GMNTu6KNYFGCvOul8+TZ1E988JaPPM8BUsnw76/Ab
%t89JnmkyHPKRN/Q9Nqpcy3p+i3s3RoqUhD6b6qgYJ/jkekGSOmcoscrw8z0TPl0TTG4TNbeHooZiRnep14vGV8sw8V659BgdrbJR
%VpvFeOXvN1imoWZerJPP7wYz+ymrhgiNcr+cX2tX0DdA9Y9kOEVEm4h6rhml+9L7NAbqSrIt89X1ullsp3a83dTxpq32ZpLBaGQK
%eJiJveWyWrK1EJfLmU9JJz19536nzXf+Sa9mPxGC7ieItawUAgGHn8qDX81O3bC0HEQ/rkGtoP68DdahBjBdUyi59TMY6YXfiNU0
%0V6FvkuwUisaZNVjG5+IRoOEBtODgzBHg+1TGal6j0gzADyTpCreOyaJ1MlnclG3QLFdZtz52WSyvd7O6ZnUxWr2ZkZHivS4hUMP
%R2y6qzzbTvLMda4X11jFrHt6c2RK3kzOuLuvdO2OziHa9fHm0srErwcFVEfEmn9Pa0SbLkeWg4rS03EZ9uOq561HfV33u6DZuvhH
%ve2yXOqsRVYyLpADUw7Pa9xjuWz8ZWbzduzJSrSQTPQYFTw6AZNISuSTnTwH8Cwn/OAH84FZGhvdHkDlDkh3xxB1jyCik808QRo8
%Na5ymPFDbFfbKzX35Sp8n6+yZ+0hDZ+eJg937E4oWPWieHpyk3y++TJwXvzvZz+9dl49+/vfz53PTSB8Mf7ofHQ+P//i/Hs6X8xo
%7Zi+ydfn8NDPky/OhcP6n+GFX3onevozDzXAxVb6iSRaVFYnwp1TDXYPbuZWD/Fu3q4W2zdvGdOP7mL9PCRBPxs7XD4pLp8MHLMy
%0ffm3BgjCa5qsNlGIRitlW1wvloZsNmfRnGOEpPZW9S+tT/t94JQMdqh/usPP7x4/vvn/1jY6MwOBuELe5yHBP/iHeCX/cv4xdoD
%6xWFwvcoxyGKWOL4v/8Alp/9/h+3n384zDPWEj16itehHfzRAZvQjgelT5jtqtqGOW4VUdPQrBKIN3Lb+0FGIN/iamdPiDr4PdTi
%xzQa7tOS99FiffkwNX6MmtynJsSfJhccI6f2yWERdz898SA9Hh2jFx1QJpfRfQSDPnvQfCKOj1GM9ykyLG/94H4hg/tJMtnzaZfn
%qD+XuSV+lBM/OiTvCUW/DMWffjfb0/e78qMpyuMUsQJWRPH56MNxkvIP5B3xMM0Xv/y+OE70DyQ72ro4SpWHRrf/9qIY30M2+gPq
%hfPfMyGUNv33H25+H9+j5PjrCZsYuId0QNvDG+e//fDj89+XX0ratq8/9E9h4PSfL5/Uz8KDOG6f40PpE1vz8B2dwzyJVNiXEbH7
%RD+5B+BwqP95DBY/wvPrT15+BuWn0J9VP1a2s7Kd6fawHB+W48NyfFiOp89Qf7KyHyv7sbJfhYfpfrLEJ0t8ssQnS3yyxEefUn+y
%sj8r+7OyPyv7V3iZ7q9K/KrEr0r8qsSvSvyqxE+fSn+ychwrx7FyHCvHsXJcRYdh3Ag6L/9ljn5G8kqfqcAA9KSkZT1jk3IifdKV
%CtOAD5MBkbEfV9w0cZ+ZG90Us7gfClGOEsrc6KYaBzWFQdhngSybyjkGyJQI+hGr4P4hEsz3lenzAI0dZGWO34c33thlQN1QqnHr
%Jh5aTcIXfYGqTDepCAyFftkURJYKeBi05AwPgRuGDvDATMd9FgIWQFb/EAu1wg6xcJCz+k7TYQoaqZDVqHVTxCw6GBRVOhB9HoYV
%EW45SKhs8pVwmgUZ2HbtV+QDjFB+2weExEBfHGSK+3HQD5l/sFFx+ITehj7EHAbFfSYDfog/gbbG62vjtfmr79DE4GM1i5Xcu/BK
%U7vwQGkTVIq2m2qt77dUptpv6asD0Kpn3HIJeQCq4iaSSGVNwIjY0t/xlr48yGN/n6wGUWbChHRV7iYgMa3NbGEG6mXeTMP1v0ag
%lCXUF8pZ9oPa7bQlYJswKknXnk+nUQYioBtwV+k84oF1e/97g0UErCGyg6PFQDdWfawJ91BZ4C4yDh0n6Zs2KmpicXnDZJ+LZkg1
%vObJJsX7TKgjDBCkptcNY4tZmxyuYwt+L7mQEO5BbRkbhsQRyYk2t+RuNcRHGho5dhqEYAcaKjOGjHioOhzhx7Yw3PUAataoCPAg
%5ofhUoV78K9935eCMIc8akPg3PIP4a9zmYrqiCmvLF+qYTbxBkopK64jr9QS8w+ZXugYaScoBbPISO5EbihQNUghxVfKsjMwgD1Z
%nZFapI5qo0r+MZKpqSNa47oR/N4Po69l7DHo2i2Y9ZmMHmC2zKwiYhpNCW5G2nQZVMMDEe0Tbmut1cSpfKtSfavBx42QsTymlh2m
%wJPY0wWC0vjvAW2I0B7eamIMpReyJj80jDUB8ShjtRTRjN3RidXSsG9DqdC19GG1cAlMQvj7LUftKq2obkItiALbgjXcYqgJ1pby
%rcCW4UEcmH8UVU9xNSVRRRyV6kb5xuA+td1bohnltYqOAK4m+M5kbAH9Y6Djk/D9o+3rbsAjVF5sD4kNx2gh6xtygWgfK4ptrBiC
%r2Y29JsJM6KiVPgHkdMahh9q4pEexuUDirIKOqop/TLNojqrSVq8MEt4ixhDEc3F44VsD7YJowV1IdunHBxRQYyqKT7UgBs4YCQe
%5ShBbCbK0o/9sh7mgkp8ETZMl6aUYWzPTiVJG6pZr4oYG80Dc0esmnnP5nmflbKBysY42lfrHjOVIoSgMBWqJthov9vSQ8shRSCr
%RQKXKDZrZ/cbBZONGnWKeGdVhnTAdyZpjoVDOc2wGO4ayv3OD0SKhUJSGbKPwLrsVnPzwfmaHR4DChVaLpEoQ6tEDKPwUbmnQcf7
%MtzDG0Q1JhurJVu3CtSDwauo+G84rxfaDZhqrIZzq8HiotVfajz7DRGh2e9+gBNa0zarBVEVUdwPmpumQrb5wAqg7sFs68laNuvS
%Bvt/4t1iva5F4z+Ptk55WO8G/u5kahG1Z9JGGSGuyim6UifkxaBqeBy24wz9GN8Ls9AXVXnCd/s9HGHVaFZFuDW6ubrfR6nyCITY
%HdNcBoeAPG5G1fjRXmeKNk2s0ZDAuKgmDR7taaWrF/J8D47sEfrC36MkxE4ENlMmq/Hb9TIsLctyV28n+dGeYRCLgdyDKmulYfEb
%HsTh7/FgTZEWCzQ3B/VWUUvbwj+gbYw7oDJ7NfCQ28Bjees2li1/4KLVysTOvXzMyqjFXFWYNo5qVbyNKuwAazyjBjastFBW/LSX
%RLZ1ZFX5c7FXxVozbV3GhgRWstl1COWhAQ/NKW00rFXrWGTtm/vnP5SqtId0aNguRmszo9UUN5s/h5KBv08IaaHZjWlTtXZYWg0H
%iOgOrX2XdqjKCLUwDw+IZvOCugaCHaAYBNZe0o5qpL1Tc7/92mywZl/ucQo5qvYuV/fwF1ibVkcZbDWIgFirSmOLbAOqdV3O3hzh
%QXqoy8UGQ8gp4fpVNWw33TdbNmKro95s8bETDkSmVomtXpSvtVvZ5rDhJGAU1fWf1RA1NpMth9vTeeM7cdXUOqdSIfQQ8rIgkHTy
%0NxiyaCa2UgfPMi6dtgZ+lCaiFDJBagmSkMJjjK78tcdVDu3D1TLWFRgXjk8Uuola1XcovxrOPha/m0yESmWy0NUeBjHLQ3uiGYr
%9LEs7OnD5iXkWCSrsGYmosqjMajF2k5TC8tuWxDcg7RlB1+1dLFjDjMXHiFCscrCg21YG/LmgGKfN3LDB3k7vk+2P4JgxMwerHaY
%BkZ56KsL7zDqC1UFeVMBtIzQVAt6FVqnsRre5shCI8rzh5IYj2yjKMxUtvKaFsJnItpOCnDvoB/JikxTvzMD/tod49Y4qSBbwMPw
%MLEHd7PpiEkEoolEe3RX0nme9MVX83g/WuZH5JHSP96BMi8XD0ohqx0OZfG5M3qHGQ6TiUD8WQnbJCXVkZE8QjGyDLZDsWXMP8xN
%yUQbc9wgblOxW9qctpraEraaGD0swFnI95uOZl1aoPlib+8anKm9Aj2gPXM/CvheILdYaop8OuHeXw8o2iLD+rrOdVRNNEIEFOV+
%tR1Ca0ymqvN6ElAFO4+hyGaF2tT9khLd13qPNYrBanWRb1F4MHRhuaqvNYx2ZSp0wtqkblY8qlmm3lPNP4zdBtv4G9Gsin1vN6GJ
%XTJCuL/gthhg0hQBeyzYS24brKxp2oIT8eAAlkOsWCtpmxN1kKDV2YJiXLNGb7MRVlOXBT66gxU3O/4NpB1vSu+v70Lam60EiR+V
%Wtq8VtHUsmu1kg7sPZsqSi01Wntitg0baNB4a3PqwZqd6Hobmx5H4e2p9f+V9y573AHRfTNc7WF2jNR4w0qftHPB5VdPZjV2+0Sm
%PkSz9lUa9K1HNGr6j3Kvduddjnc3i3aUF5IbfQ0Vy4X7+08p9IO9Gak5bao92tZKvbFU8dHArIOTuiqvHThsRgr+L/PjLi0YouAr
%XCoiw9XJlPbHyrNsOh9U9dm5jZdxqqj8R+wCtVmjPQIhxVF+Gd3UR+kVVrvH/bvPZL0D27dC7MBoJ8Uvz9N2sAfVGaVAIVQvoRv4
%wWjeo0pHSLXqaigtB8o9ll2ZSh+1tN/lUT+uztKbSxv89dmjLUXDeBNZ6k/ltZgeDdh7ZoBOEkO5l+sb8gcihtHJZuWWTHetqmfS
%XKVY9v/ZEwV05B9aCf9eOf7vPmDQ5eHXKvM/6/3byf+fP/n3/eOn/qzhpXk0LzAP8bWOI62ODz0b0wxvziNtOs1ll1w03p8WLDAc
%GeVmUE7Q9tg2yuZgYJfR7s6g5tBivye3DxVYbO/0W34VHIY3s8EuyUofj9fif+67skzYckl4K6+mWbtBtg44rBbaz249KdscYmBM
%s4lvtxwPwcbgLa4OPMlqOaE1bfHWoUxzbkC8NIcwVkNYRbANZMwmSBG1Z/J6KtwLuCiO+zEX/l7IxbxiMrR6tPs/HHk1FtqJDw8i
%ad3cf4aAhVyFsD3KokR7YqIpEejJtEg85ojYxs58xFTADsjBY5Q+0SHsko6Qgkg+qryUccNkC4vNhegrtc8BPIlH++D4MMO0rqpx
%2/0PsiJsK9mc2HALC8pQuQ890lnXvvucHF1mNV30A6S8PpgRca2ApmF3LVmCaXPAeqS2wUMFptpH8whu0L1mzJZONepiB8EtVdsN
%PD7YX9LGRdVfVs7A6XGesNrlqVVKAoWi/Ncu1Q8CmH8pqX8EQP+DF/LrMvmUv6xB/6rlX7QTwb9m/WAth+rNB76LCmqs/13H4x6c
%q9H1mx0lXj2vVblqja0pdg4SqIF0cljvgnELdbALEn2xi9riKdgHiV1cjQ6ih8nVD7PyfUQHyPG+Okzuq/7tRhDuAATfBcR/6tk3
%i/HSWW39tx7LEXtrwd1/F1LpqjE665vT9TIGtuNjMUBmt2SjEc/K7kPm+d/+vv392T/zlVz6+ptHffv7857EvOr929+3vz/uSe0/
%duT629+3v6/1JHbw+tvft7/7PImK73FZfMde6JXfgUNf/1N9acMw9GKPwa3gV8JjocfQSXks9rjvcXTnHg88Ljweelx6XHk88njs
%Bb4XhF4gvUB5gnsi8AT6UBtaeICBAY3E0ABwdOXoyANqB2qBdoziGBaEwBFGeAGpBFbJ8OJ4oU0KvIgM+kiFF/optCu0K7QptCni
%AW2K2oAjQluEsRHaI8AjwCPAY+COAY8BjzEuRluMtjj2hO/jxfCCJD5E8QVeIV7SE9CMYIAzwFmEV0y/6eiJCG3ALYAbSze80C9G
%P9AQoCFAQ8ToAzoCdATohKARgkYIGiFohKARgkboow1GCEErhBlC2CGEIUJYIgTNEPRCGCOENUKYI4Q9QhgkhEVCmCTk1Ac4oOsQ
%ug5hlxA6DqHjEDoOYZtQEAx9BfoK9A2BM0SfEH1C9AnRJ8RY2CMM0Q82CGGDEDYIJfor9IfuQ8ghoS8JWSRkkZBFQhYJWSRkkb7C
%K8IL/SCThEwS+pPQn4RsErJJyCYhm4RsEvqUkE9CPgn5JOSTkE9CPgn5JOSS8CkJn5KQT0I++pUCCX+S8CcJf5IC7ZBXQl4JeSXk
%lZBTQk4JOSVkk5BNQjYJ2ST8TcLfJPxNwt8kZJWQVUJWCVklZJWQVcK3JHxLwrckfEtGgEeAw8ckfEzCx2QEvLC/hP0l7C9hbwl7
%S+hJQU8KelLQk4KeFPSkoCcFPSnoSUFPCnpS0JOCnhT0pKAnBT0p6ElBTwp6UtCTgn4U9KOgHwW9KNhdcYJjPPSjoB8F3SjoRkE3
%CvpQ0IeCPhT0oQTaoAsFmyvYXEEPCnpQ0IOCHhTiTCHOFOJMIb6UwjjEGH0RVUQ/+unTdYRX7NHvBUTgNwK/EfiLwF8EO0awYwS7
%ReAvAn8ReIjAQyToR0MxHjxEQo6+mB95Wq4Wt3mRFhPrC8uK8psCl9s1fbPZ9az8MrV8kkUqy4NARNlkDAWmSo7HjE+zHFrJVIqG
%1DePtWkc9Y+70C90rOmHMTSeIt1u7hzei3v+ga6rfLlYz3B1Z3q/3WyW60G//2a2ebsd0zdQ9+f0K5pFX+M5gMHmOYhZmEUxS6cR
%gwtHYZCORRRO6Qv2Ix6lAeNhMMniBo35EkP98yQGxd9WGbh3RK8XO9tildPPT9FPeJZfWX3u6NRNvzJEX6FY/TCn6tUozQ+WXm03
%0+hK/0Yq4a2/vc4QNF/ZWn2RXG9ZCh/LaZr7MopVxKbpJICu2WQq0mAaRhM2jXGbq5SngmdxnrE0VBE+pQrDmIuwNoUlWPWLpgY/
%vFaMJzxOwyhD9IvIZ1Nf5CKP/WmaTqfR2J+kMoryOI0ylY05mwopJqnKUwQlmxzB3/q5MENJxGM5VhHxPc1ZAMGCaRAqyTIWy3Ts
%j5WKkJ9kFuRjGU24nLBJFPnjGF1Y+fxKSelRPwVTyjcRSHJqHENHEqE5nQZMwPAsi2WeT6KA5UGmIkyGqT+JfUxOk3DCMhFnMs8m
%wX3yHf4G3fL7BYWfT2AkhhkVkyVHkvWFHDMVZxMZ5lka5MFYIL1MgiidhhM1ngRxFjCZ5bmfhv49ZI99SWwpr0pzzNxwdzWFhCKN
%kHP4WNGzPnGMvDT1+VRFuRinfsjRS+KVxmwcTzmDou8h3Po2T0MtzKfZWEyijIVpxMeY6ujwcixzJKMUkwLCLIMGpjmJFoupmKaM
%KxVOslRN/Xi8R+2h71M0ZGFJPkmzbJJh0ghQQLBgytMs4lOpYLOMc0xS4zBPBVwuncRQeRZlmYQPw+kmvP4iR0qBX/4P0oAZ+w==
%END_EMBEDDED_REPRO_BUNDLE
